%\pdfoutput=1
%\RequirePackage[l2tabu, orthodox]{nag}

\documentclass[reqno]{amsart}

% \usepackage[notref,notcite]{showkeys} % USED TO SHOW EQUATION LABELS. REMOVE BEFORE FINAL VERSION
% \usepackage{fullpage} % FOR WIDE MARGINS

\usepackage{color,soul}


\usepackage{lmodern}
\usepackage[T1]{fontenc}
\usepackage[utf8]{inputenc}
\usepackage[english]{babel}
\usepackage{microtype} %Better font spacing

\usepackage{amsmath,amssymb,amsthm,mathrsfs,cases,
latexsym,mathtools,mathdots,booktabs,xparse,enumitem,tikz,bm,url}
%\usepackage[centertableaux]{ytableau}

\usetikzlibrary{arrows,positioning,shapes.geometric}
\usepackage[pdftex]{hyperref}
\usepackage[capitalize,noabbrev]{cleveref}

% 
% \usepackage{todonotes}
% \presetkeys%
%     {todonotes}%
%     {inline,backgroundcolor=yellow}{}

\newtheorem{theorem}{Theorem}
\newtheorem{proposition}[theorem]{Proposition}
\newtheorem{lemma}[theorem]{Lemma}
\newtheorem{corollary}[theorem]{Corollary}
\newtheorem{conjecture}[theorem]{Conjecture}

\theoremstyle{definition}
\newtheorem{definition}[theorem]{Definition}
\newtheorem{example}[theorem]{Example}
\newtheorem{remark}[theorem]{Remark}
\newtheorem{question}[theorem]{Question}
\newtheorem{problem}[theorem]{Problem}

% To make circled entries
\newcommand*\circled[1]{\tikz[baseline=(char.base)]{\node[shape=circle,draw,inner sep=2pt] (char) {#1};}}

\definecolor{lightblue}{rgb}{0.8,0.8,1.0}
\definecolor{lightgreen}{rgb}{0.8,1.0,0.8}
% Nice pastel colors
\definecolor{pBlue}{RGB}{86,139,190}
\definecolor{pCyan}{RGB}{149,186,201}
\definecolor{pSand}{RGB}{184,166,121}
\definecolor{pAlgae}{RGB}{87,115,135}
\definecolor{pSkin}{RGB}{236,216,167}
\definecolor{pGray}{RGB}{156,175,156}
\definecolor{pPink}{RGB}{215,114,127}
\definecolor{pOrange}{RGB}{211,153,80}


\newcommand{\Per}[1]{\todo[size=\tiny,inline,color=pOrange]{#1 \\ \hfill --- Per}}
\newcommand{\Frether}[1]{\todo[size=\tiny,inline,color=pBlue]{#1 \\ \hfill --- Frether}}


% The q is now included by default. Optional first parameter replaces the q with something else.
% Optional fourth parameter adds a power. In \qbinom{a}{b}[c], c is interpreted as the power.
\NewDocumentCommand\qbinom{O{q} m m o}{
\ensuremath{\IfNoValueTF{#4}{
{\genfrac{[}{]}{0pt}{}{#2}{#3}_{#1}}
}{
{\genfrac{[}{]}{0pt}{}{#2}{#3}_{#1}^{#4}}
}}}


\newcommand{\stirlingOne}[2]{{S_1(#1,#2)}}

%For defining stuff
\newcommand{\defin}[1]{%
\relax\ifmmode%
\textcolor{blue}{#1}%
\else\textcolor{blue}{\emph{#1}}%
\fi%
}


\newcommand{\oeis}[1]{\href{http://oeis.org/#1}{#1}}
\newcommand{\thsup}{\textnormal{th}}


\newcommand{\setN}{\mathbb{N}}

\newcommand{\xvec}{\mathbf{x}}
\newcommand{\yvec}{\mathbf{y}}
\newcommand{\wvec}{\mathbf{w}}



%Permutations
\newcommand{\symS}{\mathfrak{S}}
\newcommand{\dArr}{\mathfrak{D}}
\newcommand{\bdArr}{\mathfrak{BD}}

% Sets 
\newcommand{\SEF}{\mathcal{F}}  % Set of subexcedant functions
\newcommand{\DSEF}{\mathcal{DF}}%Elements in F, which are derangements



% The special involution on subexcedant functions
\newcommand{\frether}{\Psi}
\newcommand{\fretherD}{\hat{\Psi}}

\newcommand{\fretherDConj}{\hat{\Phi}}

\newcommand{\caseOne}{\ensuremath{\heartsuit}}
\newcommand{\caseTwo}{\ensuremath{\spadesuit}}
\newcommand{\caseThree}{\ensuremath{\diamondsuit}}
\newcommand{\caseFour}{\ensuremath{\clubsuit}}

\newcommand{\img}{\mathbf{m}}

%Math operators
\DeclareMathOperator{\id}{id}
\DeclareMathOperator{\inv}{inv}
\DeclareMathOperator{\sgn}{sgn}
\DeclareMathOperator{\type}{type}
\DeclareMathOperator{\length}{\ell}

\DeclareMathOperator{\des}{des}
\DeclareMathOperator{\DES}{DES}
\DeclareMathOperator{\fix}{fix}
\DeclareMathOperator{\FIX}{FIX}

\DeclareMathOperator{\exc}{exc}
\DeclareMathOperator{\EXCi}{EXCi}
\DeclareMathOperator{\EXCv}{EXCv}

\DeclareMathOperator{\rlm}{rlm}
\DeclareMathOperator{\RLMv}{RLMv}
\DeclareMathOperator{\RLMi}{RLMi}

\DeclareMathOperator{\zig}{zig}
\DeclareMathOperator{\ZIG}{ZIGx}
\DeclareMathOperator{\ZIGv}{ZIGv}
\DeclareMathOperator{\ZIGi}{ZIGi} %Same as peaks
\DeclareMathOperator{\peaks}{peaks}
\DeclareMathOperator{\PEAKS}{PEAKS}


\DeclareMathOperator{\aexc}{aexc} % strict anti-excedance
\DeclareMathOperator{\supp}{IM}

\DeclareMathOperator{\fixMap}{\mathtt{fix}}
\DeclareMathOperator{\unfixMap}{\mathtt{unfix}}
\DeclareMathOperator{\sefToPerm}{\mathtt{sefToPerm}}
\DeclareMathOperator{\flipMap}{\mathtt{flip}}
\DeclareMathOperator{\zetaMap}{\zeta}


\title{An involution on derangements preserving excedances and right-to-left minima}

\author[Per Alexandersson]{Per Alexandersson}
\address{Department of Mathematics, Stockholm University, SE-106 91 Stockholm, Sweden}
\email{per.w.alexandersson@gmail.com}

\author[Frether Getachew]{Frether Getachew}
\address{Department of Mathematics, Addis Ababa University, Ethiopia}
\email{frigetach@gmail.com}

\begin{document}

% 05A05 ; 05A19

\begin{abstract}
 We give a bijective proof of a result by R.~Mantaci and F.~Rakotondrajao from 2003,
 regarding even and odd derangement with a fixed number of excedances.
 We refine this result by also considering the set of right-to-left minima.
\end{abstract}

\maketitle 

\section{Introduction}

In this paper, we refine a theorem by R.~Mantaci and F.~Rakotondrajao from 2003. 
In order to state our results, we need to recall some standard notions and terminology.
%
Let $\defin{[n]} \coloneqq \{1,2,\dotsc,n\}$ and let $\defin{\symS_n}$
be the symmetric group acting on $[n]$.
Let $\defin{\dArr_n} \subseteq \symS_n$ denote the set of derangement
and let $\defin{\dArr^e_n}$ and $\defin{\dArr^o_n}$ be the subsets
of even and odd derangement, respectively.
%
For any function $g:[n]\longrightarrow [n]$, let 
the set of \defin{exedances}, \defin{exedance values},
\defin{right-to-left minima indices}, and \defin{right-to-left minima values}
respectively, be defined as 
\begin{align*}
\defin{\EXCi(g)} &\coloneqq \{ j \in [n] : g(j)>j \}, \\
\defin{\EXCv(g)} &\coloneqq \{ g(j) : j \in \EXCi(g) \}, \\
\defin{\exc(g)} &\coloneqq |\EXCi(g)| = |\EXCv(g)|, \\
\defin{\RLMi(g)} &\coloneqq  \{i\in [n]: g(i)<g(j) \text{ for all } j\in \{i+1,\dotsc,n\} \},\\
\defin{\RLMv(g)}  &\coloneqq  \{ g(i) : i \in \RLMi(g)\}, \\
\defin{\rlm(g)} &\coloneqq |\RLMi(g)| = |\RLMv(g)| .
\end{align*}
More details and missing definitions are provided later in \cref{sec:prelim}.

It was shown by R.~Mantaci and F.~Rakotondrajao~\cite[Cor. 5.11]{MantaciRakotondrajao2003},
that for every $n\geq 1$ and $1 \leq k \leq n-1\}$,
\begin{equation}\label{eq:mainEquation}
 |\{ \pi \in \dArr^e_n  : \exc(\pi) = k \}|
 -
 |\{ \pi \in \dArr^o_n  : \exc(\pi)=k \}|
 =(-1)^{n-1}.
\end{equation}
This refines a result by Chapman's, stating that
$|\dArr^e_n|-|\dArr^o_n|=(-1)^{n-1}(n-1)$, see \cite{Chapman2001}.

%
The main result of our article is a
proof of a refinement of \eqref{eq:mainEquation}.
In \cref{thm:theMainTheorem}, we prove that
\begin{equation}\label{eq:mainResult}
 \sum_{\pi \in \dArr_n} (-1)^{\inv(\pi)} \!
 \left( \hspace{5mm}
 \prod_{\mathclap{j\in \RLMv(\pi)}}  x_j 
 \hspace{3mm}
 \right)
 \!\!
 \left( \hspace{5mm}
 \prod_{\mathclap{j\in \EXCv(\pi)}}  y_{j}
 \hspace{3mm}
 \right)
 = 
 (-1)^{n-1} \sum_{j=1}^{n-1} x_1 \!\dotsm \! x_j y_{j+1}  \!\dotsm  \! y_{n}.
\end{equation}
%
We prove \eqref{eq:mainResult} by
exhibiting a bijection $\frether: \dArr_n \to \dArr_n$
with exactly $(n-1)$ fixed-points,
and where $\frether$ acts as a sign-reversing involution
outside the set of fixed-points.
The bijection preserves all
the relevant permutation statistics,
which gives the desired result.
Moreover, \eqref{eq:mainResult} allows us to deduce that
\begin{equation}\label{eq:mainResult2}
\sum_{\pi \in \dArr_n} (-1)^{\inv(\pi)} \!
 \left( \hspace{5mm}
 \prod_{\mathclap{j\in \RLMi(\pi)}}  x_j 
 \hspace{3mm}
 \right)
 \!\!
 \left( \hspace{5mm}
 \prod_{\mathclap{j\in \EXCi(\pi)}}  y_{j}
 \hspace{3mm}
 \right)
 = 
 (-1)^{n-1} \sum_{j=1}^{n-1} y_1 \!\dotsm \! y_j x_{j+1}  \!\dotsm  \! x_{n}.
\end{equation}
where we now consider indices instead of values.

\bigskip 

The following theorem was proved by R.~Mantaci, albeit stated in a slightly different manner.
\begin{theorem}[{See  R.~Mantaci~\cite{Mantaci1993}}]
Let $n \geq 1$, then
\begin{equation}\label{eq:robertoExc}
\sum_{\substack{\pi \in \symS_n}} (-1)^{\inv(\pi)} 
 \left(\prod_{j \in \EXCi(\pi)} x_j \right)
= \prod_{j \in [n-1]} (1-x_j).
\end{equation}
\end{theorem}
We include a proof of \eqref{eq:robertoExc}
in \cref{sec:genFunc}, as it is used in \cref{thm:mainGenFunc},
where we give a short generating function proof of the $y_1=y_2=\dotsb = y_n$
case of \eqref{eq:mainResult}.
%
In \cref{sec:genFunc} we prove the following
right-to-left-minima analog of \eqref{eq:robertoExc}.
\begin{theorem}
Let $n \geq 1$, then
\begin{equation}\label{eq:rlminSumIntro}
 \sum_{\substack{\pi \in \symS_n   }} (-1)^{\inv(\pi)} 
 \left(\prod_{j \in \RLMv(\pi)} x_j \right)
 = 
	\big(\prod_{\substack{i \in [n] \\ i \text{ odd}}} x_i\big)
	\big(\prod_{\substack{j \in [n] \\ j \text{ even}}} (x_j-1)\big).
\end{equation}
\end{theorem}



\section{Preliminaries}\label{sec:prelim}


The number of \defin{inversions} of a function $g:[n]\to [n]$ is
\begin{equation}\label{eq:invDef}
\defin{\inv(g)} \coloneqq |\{ (i,j) : 1\leq i < j \leq n \text{ such that } g(i)>g(j) \}|.
\end{equation}
%
Recall that $\pi$ is \defin{even} if $\inv(\pi)$ is even, otherwise $\pi$ is called \defin{odd}. 
The set of even permutations in $\symS_n$ is denoted $\defin{\symS^e_n}$,
and the set of odd permutations is $\defin{\symS^o_n}$.

\medskip 

The \defin{fixed point set} of a permutation $\pi \in \symS_n$ is defined as
$
  \defin{\FIX(\pi)} \coloneqq \{ i \in [n] : \pi(i) = i \},
$
and we set $\defin{\fix(\pi)} \coloneqq |\FIX(\pi)|$.
Recall that the set of derangements is defined as
$
 \defin{\dArr_n} \coloneqq \{ \pi \in \symS_n : \fix(\pi) = 0 \},
$
and that $\defin{\dArr^o_n}$ and $\defin{\dArr^e_n}$
are the sets of odd and even derangements, respectively, in $\dArr_n$.

Whenever $S = \{s_1,\dotsc,s_m\}$ is a finite set of positive integers,
we shall let $\defin{\xvec_S}$ denote the product $x_{s_1}x_{s_2}\dotsm x_{s_m}$.
By definition, $\defin{\xvec_\emptyset} \coloneqq 1$.


\begin{example}
Below we show three permutations in $\symS_7$.
The first permutation has $3$ and $6$ as fixed-points, so it is not a derangements,
while the remaining two are.
\begin{center}
 \begin{tabular}{l c llll}
  Permutation, $\pi$ & $\inv(\pi)$ & $\EXCi(\pi)$ &  $\RLMi(\pi)$ & $\RLMv(\pi)$  \\
  \midrule 
 2135764 & 5 & 145 & 237 & 134  \\
 2153746 & 5 & 135 & 2467 & 1346  \\
 6713245 & 11 & 12 & 3567 & 1245 
 \end{tabular}
\end{center}
\end{example}

The right-to-left minima statistic and the excedance statistic 
behave quite differently. It is classically known that
\[
   \sum_{\pi \in \symS_n} t^{\rlm(\pi)} = 
   \sum_{\pi \in \symS_n} t^{\length(\type(\pi))} =
   \sum_{k=1}^n \stirlingOne{n}{k} t^k 
\]
where $\stirlingOne{n}{k}$ is the unsigned Stirling number of the first kind,
see \oeis{A008275}.
%
However,
\[
 \sum_{\pi \in \symS_n} t^{\exc(\pi)} = 
 \sum_{k=1}^{n} A_{n,k} t^{k-1} 
\]
where $A_{n,k}$ denote the Eulerian numbers, \oeis{A008292}.


\subsection{Subexcedant functions}

The involution we shall construct is not performed directly on permutations,
but rather on so-called \emph{subexcedant functions} which are in bijection with permutations.
Our main reference is \cite{MantaciRakotondrajao2001},
were several fundamental properties are proved.
%
\begin{definition}\label{subex}
	A \defin{subexcedant function}  $f$ on $[n]$ is a map $f : [n] \longrightarrow [n]$
	such that
	\begin{equation*}
		1\leq f(i)\leq i \text{ for all } 1\leq i\leq n.
	\end{equation*}
	We let $\defin{\SEF_n}$ denote the set of all 
	subexcedant functions on $[n]$. 
	The \defin{image} of $f \in \SEF_n$ is defined as
	$\defin{\supp(f)} \coloneqq \{ f(i) : i \in [n] \}$.
\end{definition}
We write subexcedant functions as words, $f(1)f(2)\dotsc f(n)$.
For example, the subexcedant function $f=112352$, have $\supp(f)=\{1,2,3,5\}$.

From each subexcedant function $f \in \SEF_{n-1}$,
one can obtain $n$ distinct subexcedant functions in $\SEF_{n}$ by appending
any integer $i\in [n]$ at the end of the word representing $f$.
Hence, the cardinality of $\mathcal{F}_n$ is $n!$.
There is a bijection $\defin{\sefToPerm}: \SEF_n \longrightarrow \symS_n$,
described in \cite{MantaciRakotondrajao2001}, 
which is defined as the following product:
\begin{equation}\label{eq:sefToPerm}
  \defin{\sefToPerm(f)} \coloneqq (1\,\,f(1)) (2\,\,f(2)) \dotsm (n\,\,f(n)).
\end{equation}

\begin{example}\label{ex:sefToPerm}
Let $f = 112435487 \in \SEF_9$.
The permutation $\sigma=\sefToPerm(f)$ is
 \begin{align*}
 	\sigma &=(1)(2\,\,1)(3\,\,2)(4)(5\,\,3)(6\,\,5)(7\,\,4)(8)(9\,\,7)\\
 	&=(1\,\,6\,\,5\,\,3\,\,2)(4\,\,9\,\,7)(8)\\
 	&=\begin{pmatrix}
 	1\,2\,3\,4\,5\,6\,7\,8\,9 \\ 6\,1\,2\,9\,3\,5\,4\,8\,7
 	\end{pmatrix}.
 \end{align*}
\end{example}




For $\sigma \in \symS_n$ and $j \in [n]$,
it is fairly straightforward to see that we can compute
the $j^\thsup$ entry of $\sefToPerm^{-1}(\sigma)$
via the recursive formula
\begin{equation}\label{eq:permToSefRecursion}
 \sefToPerm^{-1}(\sigma)_j \coloneqq 
\begin{cases}
 \sigma(n) \text{ if $j=n$,} \\
 \sefToPerm^{-1}( \; \sigma \circ (n\,\,\sigma(n)) \;  )_j &\text{ otherwise}.
\end{cases}
\end{equation}
Note that 
\begin{equation}\label{eq:permProjection}
 \sigma' \coloneqq \sigma \circ (n\,\,\sigma(n))
\end{equation}
is the result after interchanging the image of $n$ and $n$.
Therefore, $\sigma'(n)= n$ and $\sigma'$ can be considered as a permutation in $\symS_{n-1}$.
Hence, the definition above is well-defined, and for simplicity,
we use the shorthand $\defin{f_\sigma} \coloneqq \sefToPerm^{-1}(\sigma)$.

\begin{example}\label{ex:permToSef}
We shall now show how to invert the calculation in \cref{ex:sefToPerm}.
We start with the permutation $\sigma^9=\binom{1\,2\,3\,4\,5\,6\,7\,8\,9}{6\,1\,2\,9\,3\,5\,4\,8\,7}$,
and for $i \geq 1$ we let $\sigma^{i-1} \in \symS_{i-1}$ be given by
\[
  \sigma^{i-1} \coloneqq \sigma^{i} \circ (i\,\,\sigma^i(i)),
\]
where we use the observation in \eqref{eq:permProjection}.
Combining this recursion with \eqref{eq:permToSefRecursion},
we have
\begin{equation*}
\begingroup
\renewcommand*{\arraystretch}{1.2}
\begin{array}{lr}
	\sigma^9=\binom{1\,2\,3\,4\,5\,6\,7\,8\,9}{6\,1\,2\,9\,3\,5\,4\,8\,\mathbf 7}  & f_{\sigma}(9)=7 \\
	\sigma^8=\binom{1\,2\,3\,4\,5\,6\,7\,8}{6\,1\,2\,7\,3\,5\,4\,\mathbf 8 } & f_{\sigma}(8)=8 \\
	\sigma^7=\binom{1\,2\,3\,4\,5\,6\,7}{6\,1\,2\,7\,3\,5\, \mathbf 4} & f_{\sigma}(7)=4 \\
	\sigma^6=\binom{1\,2\,3\,4\,5\,6}{6\,1\,2\,4\,3\,\mathbf  5} & f_{\sigma}(6)=5 \\
	\sigma^5=\binom{1\,2\,3\,4\,5}{5\,1\,2\,4\,\mathbf  3} & f_{\sigma}(5)=3 \\
	\sigma^4=\binom{1\,2\,3\,4}{3\,1\,2\,\mathbf  4} & f_{\sigma}(4)=4 \\
	\sigma^3=\binom{1\,2\,3}{3\,1\,\mathbf  2} & f_{\sigma}(3)=2 \\
	\sigma^2=\binom{1\,2}{2\,\mathbf  1} & f_{\sigma}(2)=1 \\
	\sigma^1=\binom{1}{\mathbf 1} & f_{\sigma}(1)=1.
\end{array}
\endgroup
\end{equation*}
Thus, $f_{\sigma}=112435487$.
\end{example}



\begin{proposition}[{See \cite[Prop. 3.5]{MantaciRakotondrajao2001}}]\label{prop:excAsIm}
For $f_{\sigma} \in \SEF_n$ we have that
$
 [n] \setminus \supp(f_{\sigma}) = \EXCv(\sigma).
$
In particular, $\exc(\sigma) = n - |\supp(f_{\sigma})|$.
\end{proposition}

Since subexcedant functions are seen as maps $g:[n] \to [n]$,
we have the notion of right-to-left minima, fixed points, etc., as defined in the introduction.
%
The following proposition is reminiscent of some results in \cite{BarilVajnovszki2017},
but they consider a different map between permutations and subexcedant functions.
%
\begin{proposition}\label{prop:rlmForSEF}
Let $\pi\in\symS_n$ and $f_{\pi}$ be the corresponding subexcedant function. 
Then
\begin{enumerate}
    \item $i \in \RLMi(\pi)$ implies $\pi(i)=f_{\pi}(i)$,
    \item $\RLMv(\pi) = \RLMv(f_{\pi})$,
    \item $\RLMi(\pi) = \RLMi(f_{\pi})$.
\end{enumerate}
\end{proposition}
\begin{proof}
We use induction over $n$; the base case for $n=1$ is trivial.
Assume now that for any $\sigma \in \symS_{n-1}$, 
\begin{enumerate}[label=(\roman*)]
    \item $i\in \RLMi(\sigma)$ implies $\sigma(i)=f_{\sigma}(i)$,
    \item $\RLMv(\sigma) = \RLMv(f_{\sigma})$,
    \item $\RLMi(\sigma) = \RLMi(f_{\sigma})$.
\end{enumerate}
%
Now let $\pi^n \in\symS_n$ and define $\pi^{n-1} \in \symS_{n-1}$  as
\begin{align}\label{eq:someLemmaRec}
    \defin{\pi^{n-1}(j)} \coloneqq \begin{cases}
             \pi^{n}(n) &\text{if }\pi^n(j)=n\\
             \pi^{n}(j) &\text{otherwise},
			\end{cases}
\; \text{ so that } \; 
f_{\pi^n}(j)=\begin{cases}
             \pi^n(n) &\text{if }j=n\\
             f_{\pi^{n-1}}(j) &\text{otherwise.}
             \end{cases}
\end{align}
This is the same setup as in \cref{ex:permToSef}.
Now suppose $i \in \RLMi(\pi)$. 
By induction hypothesis, we have that 
$\pi^{n-1}(i) = f_{\pi^{n-1}}(i)$ if 
$\pi^{n-1}(i) \in \RLMv(\pi^{n-1})$, $\RLMv(\pi^{n-1})=\RLMv(f_{\pi^{n-1}})$,
and $\RLMi(\pi^{n-1}) =  \RLMi(f_{\pi^{n-1}})$.


\medskip 
\noindent
\textbf{Proof of property (1).}
We must show that $\pi^n(i)=f_{\pi^n}(i)$.

\noindent
\textbf{Case $i=n$}:
 Here, $\pi^n(i) = f_{\pi^n}(i)$, as this follows immediately \eqref{eq:permToSefRecursion}.

\noindent
\textbf{Case $i<n$}: 
Now, either $\pi^n(i)=n$ or $\pi^n(i)=\pi^{n-1}(i)$,
since for any $j \in [n]$,
\begin{equation*}
    \pi^n(j)=\begin{cases}
             \pi^n(n) &\text{ if }j=n,\\
             n &\text{ if }j<n\text{ and } \pi^{n-1}(j)=\pi^n(n),\\
             \pi^{n-1}(j) &\text{ otherwise.}
             \end{cases}
\end{equation*}
But $\pi^n(i)=n$ is impossible since $\pi^n(i)$ is a right-to-left minima and $i<n$.
Hence 
\begin{equation}\label{eq:someLemma}
 \pi^n(i)=\pi^{n-1}(i) \text{ and } f_{\pi^n}(i)=f_{\pi^{n-1}}(i).
\end{equation}
\medskip 
Since $i$ is a right-to-left-minima of $\pi^n$ we have
\[
 \pi^n(i) < \pi^n(j) \text{ whenever $i < j \leq n$}.
\]
By the first formula in \eqref{eq:someLemmaRec}, 
we have that $\pi^{n-1}(j)=\pi^{n}(j)$ whenever $i<j<n$
(because $\pi^n(i) \neq n$).
Hence,
\[
 \pi^{n-1}(i) = \pi^n(i) < \pi^n(j) = \pi^{n-1}(j) \text{ whenever } i < j \leq n-1,
\]
so $i \in \RLMi(\pi^{n-1})$. This fact together
with \eqref{eq:someLemma} and the induction hypothesis, 
finally gives
\[
i\in \RLMi(\pi^n) \implies \pi^n(i) = \pi^{n-1}(i)=f_{\pi^{n-1}}(i)=f_{\pi^n}(i).
\]

\medskip 
\noindent
\textbf{Proof of property (2).}
By definition of $\pi^{n-1}$, we have that
\begin{align*}
    \RLMv(\pi^n) &=(\RLMv(\pi^{n-1})\cap [\pi^n(n)]) \; \cup \; \{\pi^n(n)\},  \\
    &= \{ \text{by induction hypothesis} \} \\
    &=(\RLMv(f_{\pi^{n-1}})\cap[f_{\pi^n}(n)])\cup\{f_{\pi^n}(n)\}\\
    &=\RLMv(f_{\pi^n}).
\end{align*}

\medskip 
\noindent
\textbf{Proof of property (3).}
By the first property together with the inductive hypothesis,
we have
\begin{align*}
    \RLMi(\pi^n) &= \{j \in \RLMi(\pi^{n-1}) : \pi^{n-1}(j) \leq \pi^n(n) \} \; \cup \; \{n\} \\
    &=\{j \in \RLMi(f_{\pi^{n-1}}): f_{\pi^{n-1}}(j) \leq f_{\pi^n}(n)\} \; \cup \; \{n\}\\
    &=\RLMi(f_{\pi^n}).
\end{align*}
This concludes the proof.
\end{proof}


We say that $f$ has a \defin{strict anti-excedance} at $i$ if $f(i)<i$.
Let $\defin{\aexc(f)}$ denote the number of strict anti-excedances in $f$.
The subexcedant function $f_\sigma$
is \defin{even} \defin{(odd)} if $\sigma$ is even (odd).
\begin{proposition}[{See \cite[Prop. 4.1]{MantaciRakotondrajao2001}}]\label{prop:evenSEF}
The permutation $\sigma$ is even if and only if $\aexc(f_{\sigma})$
is even.
\end{proposition}


A \defin{fixed point} of $f \in \SEF_n$ is an integer $i \in [n]$ such that $f(i)=i$.
Moreover, $i$ is a \defin{multiple fixed point} of $f$ if:
\begin{enumerate}
	\item $f(i)=i$ and 
	\item there is some $j>i$ such that $f(j)=i$.
\end{enumerate}

\begin{proposition}[{See \cite[Prop. 3.8]{MantaciRakotondrajao2001}}]\label{prop:derangmentSEF}
We have that $\sigma \in \dArr_n$ if and only if all
fixed points of $f_{\sigma}$ are multiple.
\end{proposition}





\section{An involution and its consequences}\label{sec:involution}


The subexcedant function $f$ is \defin{matchless}
if it is of the form
\[
\defin{ f_{\sigma_{n,k}} } \coloneqq 1\,1\,2\,3\,4\dotsc k{-}1\,\,k\,\,k\dotsc k \quad \text{  for $1\leq k \leq n{-}1$}.
\]
There are $n-1$ matchless subexcedant functions of length $n$.
For example, for $n=10$, the following subexcedant functions are matchless:
\begin{align*}
 &1111111111, \quad 1122222222, \quad 1123333333, \\
 &1123444444, \quad 1123455555, \quad 1123456666, \\
 &1123456777, \quad 1123456788, \quad 1123456789. 
\end{align*}


\begin{lemma}[Properties of matchless functions]\label{lem:propertiesOfMatchless}
Let $f_{\sigma_{n,k}} \in \DSEF_n$ be matchless.
Then 
 \begin{align*}
  \sigma_{n,k} &=
(1\,\,k{+}1\,\,k{+}2 \; \dotsc\; n\,\,k\,\,k{-}1 \; \dotsc \; 2), \\
(-1)^{\inv(\sigma_{n,k})} &= (-1)^{n-1}, \\
\EXCv(\sigma_{n,k}) &= [n]\setminus [k],\text{ and }\\
\RLMv(\sigma_{n,k}) &= [k].
\end{align*}
\end{lemma}
\begin{proof}
The form of $\sigma_{n,k}$ follows directly from \eqref{eq:sefToPerm}.
Since $\sigma_{n,k}$ has only one cycle, its sign is $(-1)^{n-1}$.

Looking directly at the definition of $f_{\sigma_{n,k}}$, 
we have that
\[
 \supp(f_{\sigma_{n,k}}) = [k] \implies \EXCv(\sigma_{n,k}) = [n]\setminus [k],
\]
by \cref{prop:excAsIm}. Similarly, the last property follows from 
\cref{prop:rlmForSEF}.
\end{proof}
Note that for each $k \in [n]$, there is a unique matcheless 
subexcedant function, with $k$ exedances.
We shall see that this property is exactly what explains 
R.~Mantaci and F.~Rakotondrajao's identity in~\cref{eq:mainEquation}.
\vspace{1cm}


Let 
\begin{equation}
 \defin{\DSEF_n} \coloneqq \{ f_{\sigma} : \sigma\in \dArr_n \}
 %\text{ and $f_{\sigma}$ is not matchless}
 \text{ and }
  \defin{\DSEF^*_n} \coloneqq \{ f_{\sigma} : \sigma\in \dArr_n 
 \text{ and $f_{\sigma}$ is not matchless}\}.
\end{equation}
Hence, $\DSEF_n$ is the set of subexcedant functions 
corresponding to dearrangements of $[n]$.
% Since all matchless subexcedant functions correspond 
% to derangements, we have that $\DSEF_n$
% has cardinality $|\dArr_n| - (n-1)$.
%
Note that every $f \in \DSEF_n$ must have at least two $1$s
in its row representation.

Our goal is now to define an involution $\fretherD:\DSEF_n\longrightarrow \DSEF_n$, 
with the following properties.
\begin{enumerate}[label=(\roman*)]
\item 
	The image is preserved, $\supp(\fretherD(f_\sigma )) = \supp( f_\sigma )$.
\item 
	The set of right-to-left minima is preserved,
	$\RLMv( \fretherD(f_\sigma ) ) = \RLMv( f_\sigma )$.
\item 
	The fixed-points of $\fretherD$ consist of the matchless subexcedant functions.
\item 
	The sign is reversed, $\sgn(\fretherD(f_\sigma )) = -\sgn( f_\sigma )$,
	whenever $f_\sigma \in \DSEF_n^*$.
\end{enumerate}

\bigskip 

We shall define $\defin{\fretherD}:\DSEF_n \longrightarrow \DSEF_n$ below,
where $f_{\tau}$ is short for $\fretherD(f_{\sigma})$.
First, if $f_\sigma$ is matchless, we set $\defin{\fretherD(f_\sigma)}\coloneqq f_\sigma$.
%
Now we fix some $f_{\sigma}\in \DSEF^*_n$ and 
let 
\begin{equation}\label{eq:suppGeq1}
  \{\defin{ \img_1, \img_2, \img_3, \dotsc, \img_{\ell} } \}
\end{equation}
be the elements in $\supp(f_{\sigma})$
ordered increasingly. Note that $\img_1=1$ and  since $f_\sigma$ is non-matchless,
we know that $\ell \geq 2$ in \eqref{eq:suppGeq1}.
%
With these preparations, we define two auxiliary maps, 
$\fixMap_i$, $\unfixMap_i$ on subexcedant functions.
For $i \in \{2,\dotsc,\ell\}$,
\begin{equation}\label{eq:auxmaps}
   \defin{\fixMap_i}(f_\sigma)(\img_i) \coloneqq \img_i, \qquad 
   \defin{\unfixMap_i}(f_\sigma)(\img_{i}) \coloneqq\img_{i-1}
\end{equation}
while the remaining entries are untouched.
%
Let the conditions
\begin{equation}\label{eq:case34Conditions}
\begin{aligned}
  f_{\sigma}(\img_i) < \img_i < \img_\ell,\quad
  f^{-1}_\sigma(1) = \{1,2\}, \text{ and }
  \{\img_i + 1\}\subsetneq f^{-1}_\sigma(\img_i)
\end{aligned} \tag{\text{$\ast$}}
\end{equation}
be referred to as \eqref{eq:case34Conditions}.
Note that 
\[
  \{\img_i + 1\}\subsetneq f^{-1}_\sigma(\img_i) \iff 
  f_\sigma(\img_i + 1) = \img_i \text{ and } |f^{-1}_\sigma(\img_i)| \geq 2.
\]

\medskip 

Now let $\img_i \geq 2$ be the \emph{smallest} element
in \eqref{eq:suppGeq1} satisfying one of the cases below,
and let $f_\tau$ be given as described in each case.
%
%
\begin{itemize}[leftmargin=2cm]
 \item[\textbf{Case \caseOne:}] If $f_{\sigma}(\img_i) = \img_i$,
	then $f_\tau \coloneqq \unfixMap_i(f_\sigma)$.

 \item[\textbf{Case \caseTwo:}] If $f_{\sigma}(\img_i) < \img_i$ and $|f^{-1}_\sigma(1)| \geq 3$,
 then $f_\tau \coloneqq \fixMap_i(f_\sigma)$.
 
 \item[\textbf{Case \caseThree:}] \label{case3a} 
If \eqref{eq:case34Conditions} holds and $f_\sigma(\img_{i+1}) = \img_{i+1}$, then
$
 f_\tau \coloneqq \unfixMap_{i+1}(f_\sigma).
$

\item[\textbf{Case \caseFour:}] \label{case4} 
If \eqref{eq:case34Conditions} holds and $f_\sigma(\img_{i+1}) < \img_{i+1}$, then
$
 f_\tau \coloneqq \fixMap_{i+1}(f_\sigma).
$
\end{itemize}


\begin{remark}\label{rem:case2}
 Suppose (\caseTwo) applies for some $i \geq 2$.
 Then, for sure $f(\img_2) < \img_2$, since otherwise, we would be in case $(\caseOne)$ for $i=2$.
 Hence, (\caseTwo) only applies for $i=2$.
\end{remark}



We have several things that needs to be proved.
In \cref{lem:fretherDIsWellDefined} we show that $\fretherD$ is well-defined,
and in \cref{lem:correctRange}, we show that the range is correct.
%
In \cref{lem:imPreserving}, we show that $\fretherD$ preserves the image.
In \cref{lem:signReversing}, we show that $\fretherD$ is sign-reversing on $\DSEF^*_n$.

\begin{lemma}[Well-defined]\label{lem:fretherDIsWellDefined}
Let $f_\sigma \in \DSEF_n$. If none of the four cases 
($\caseOne,\caseTwo,\caseThree,\caseFour$) above
apples to $f_\sigma$, then $f_\sigma$ is matchless.

Moreover, if no $i \in \{2,\dotsc,t\}$ fulfills any of the four
conditions for some $t\geq 2$ and for $i=t+1$, \eqref{eq:case34Conditions} does not hold,
then 
\begin{equation}\label{eq:sefPartial}
 f_\sigma(1)\; f_\sigma(2)\; \dotsc \; f_\sigma(t)  \; f_\sigma(t+1) 
\end{equation}
is matchless. Also, if the last two values in \eqref{eq:sefPartial} coincide,
then $f_\sigma(t+1)  = \img_\ell = \ell$ and $f_\sigma$ is the matchless subexcedant function
\begin{equation}\label{eq:sefMatchless}
 f_\sigma = 1\; 1\;2\;3\; \dotsc \; \ell{-}1 \; \ell \; \ell \; \dotsc \; \; \ell.
\end{equation}
\end{lemma}
%
\begin{proof}
We first note that if $\ell=1$, then $f_{\sigma}=1\,1\,1\,\dotsc \,1$ (which is matchless)
so from now, assume $\ell\geq 2$. 
We shall use induction on $t$.

\textbf{Base case $t=2$}: Suppose that none of the four cases holds 
for $i=2$, and that \eqref{eq:case34Conditions} does not hold for $i=3$.

If Case {\caseThree} is not fulfilled for $i=2$, then
$f_{\sigma}(\img_2)<\img_2$ so $f_{\sigma}(\img_2)=1$.
Hence, $f_{\sigma}(1)=1$ and 
$f_{\sigma}(2)=1$, otherwise $f_{\sigma}(2)=2$ 
which would violate our assumption.

If Case {\caseFour} is not fulfilled for $i=2$, then
$f_{\sigma}(\img_2)<\img_2$ and $|f_{\sigma}^{-1}(1)|<3$.
Thus, $f_{\sigma}^{-1}(1)=\{1,2\}$.

If Case {\caseThree} is not fulfilled for $i=3$, then
$f_{\sigma}(3)<3$ and it follows that $f_{\sigma}(3)=2$.
Thus, $f_{\sigma}$ starts as $1\,1\,2$ so \eqref{eq:sefPartial}
is matchless for $t=2$.

\bigskip 
\noindent
\emph{Induction hypothesis:}
Suppose the statements hold for $t=k$, 
for some $k>2$. We shall prove that they hold for $t=k+1$.

Suppose that none of the cases (\caseOne,\caseTwo,\caseThree,\caseFour)
hold for $i \in \{2,3,\dotsc,k+1\}$,
and for $i=k+2=t+1$, \eqref{eq:case34Conditions} is not satisfied.

By the induction hypothesis, either
\begin{align*}
 &\ell \in [k-1] \quad \text{ and } \quad 
 &&f_\sigma(1)\; f_\sigma(2)\; \dotsc \; f_\sigma(n) = 1\;1\;2\;3\;\dotsc\;\ell\; \ell\;\dotsc\;\ell
 \quad  \text{ or} \\
 & \ell \geq k \quad \text{ and } \quad  &&f_\sigma(1)\; f_\sigma(2)\; \dotsc \; f_\sigma(k+1) 
 = 1\;1\;2\;3\; \dotsc\; k{-}1 \; k.
\end{align*}
   
We want to show that either $\ell\in [k]$ and 
   $f_{\sigma}=1\,1\,2\,3\,\dotsb\,\ell\,\ell\,\dotsb\,\ell
     $ or $f_{\sigma}$ starts with 
   $1\,1\,2\,3\,\dotsb\,k\,\,k{+}1$.
If $\ell\in [k-1]$, then we are done by the induction hypothesis.

If $\ell \geq k$, then $f_{\sigma}$ starts with $1\,1\,2\,3\,\dotsb\,k$.
Now, the possible values of $f_{\sigma}(k+2)$ are in $[k+2]$.
%
\begin{itemize}
 \item 
If $f_{\sigma}(k+2)=1$, then we are in case (\caseTwo) for $i=2$, violating our assumption.

\item 
If $f_{\sigma}(k+2)\in \{2,3,\dotsc,k\}$, then we 
satisfy \eqref{eq:case34Conditions} for $i=f_{\sigma}(k+2)$,
unless $i=\ell$ since $f_{\sigma}(\img_i)=\img_{i-1}<\img_i$ and 
$\{\img_i+1\}\subsetneq f_{\sigma}^{-1}(\img_i)$. Since 
$i\leq k\leq \ell$, it follows that $i=k=\ell$, $f_\sigma(k+2)=k$ and
and by induction, the statement in \eqref{eq:sefMatchless} tells us that
$f_{\sigma}=1\;1\;2\;3\;\dotsc\;\ell\; \ell \;\dotsc\; \ell$.

\item 
If $f_{\sigma}(k+2)=k+1$, then $f_{\sigma}$ 
starts with $1\,1\,2\,3\,\cdots\,k\,\,k{+}1$ so we are done.

\item 
If $f_{\sigma}(k+2)=k+2$, then we are in 
case (\caseOne) for $i=k+1$ or $i=k+2$, depending whether $k+1 \in 
\supp(f_{\sigma})$, again violating the assumption.
\end{itemize}
\end{proof}




\begin{lemma}[Correct range]\label{lem:correctRange}
If $f_{\sigma}\in \DSEF^*_n$, then $\fretherD(f_{\sigma})\in \DSEF^*_n$.
\end{lemma}
\begin{proof}
Let $\img_i$ be the smallest element in 
$\supp(f_{\sigma})=\{\img_1,\img_2,\ldots,\img_\ell\}$ 
satisfying one of the cases in (\caseOne,\caseTwo,\caseThree,\caseFour),
and set $f_\tau \coloneqq \fretherD(f_{\sigma})$.
By \cref{prop:derangmentSEF}, it suffices to show that all fixed-point of $f_\tau$ are multiple.

If $\img_i$ fulfills case (\caseOne), then 
$f_{\tau}(\img_i)=\img_{i-1}$. Since $\img_i$ is the 
smallest element with the property $f_{\sigma}(\img_i)=\img_i$, 
then $\img_{i-1}$ is not a fixed point in both 
$f_{\sigma}$ and $f_{\tau}$.

If $\img_i$ fulfills case (\caseTwo), then 
$f_{\tau}(\img_i)=\img_{i}$ and there is some $j>\img_i$ such 
that $f_{\tau}(j)=f_{\sigma}(j)=\img_{i}$. That is, 
$\img_i$ is a multiple fixed point in $f_{\tau}$.

If $\img_i$ fulfills case (\caseThree), then $f_{\sigma}(\img_{i+1})=\img_{i+1}$
and neither $\img_{i}$ nor $\img_{i+1}$ are fixed point for $f_{\tau}$,
since 
\[
 f_{\tau}(\img_{i+1}) = \img_{i} > f_{\sigma}(\img_i) = f_{\tau}(\img_i),
\]
due to the conditions in \eqref{eq:case34Conditions}.

If $\img_i$ fulfills case (\caseFour), $f_{\sigma}(\img_{i+1})<\img_{i+1}$
and there is some $j>\img_{i+1}$ such that $f_{\tau}(\img_{j})=f_{\sigma}(j)=\img_{i+1}$.
Consequently, $f_{\tau}(\img_{i+1})=\img_{i+1}$ is a \emph{multiple} fixed point in $f_{\tau}$,
so $\tau$ is also a derangement.
\end{proof}

\begin{lemma}[Image-set preserving]\label{lem:imPreserving}
For $f_{\tau}=\fretherD(f_{\sigma})$, we have 
\begin{equation}\label{eq:imPreserving}
\supp(f_{\sigma}) = \supp(f_{\tau}) \text{ and }
\EXCv(\sigma) = \EXCv(\tau). 
\end{equation}
\end{lemma}
\begin{proof}
Let $\img_i$ be the smallest element in 
$\supp(f_{\sigma})=\{\img_1,\img_2,\dotsc,\img_\ell\}$ 
satisfying one of the cases in (\caseOne,\caseTwo,\caseThree, \caseFour).

If $\img_i$ fulfills case (\caseOne), then 
$f_{\sigma}(\img_i)=\img_i$ and there is some $j>\img_i$ such 
that $j\in f_{\sigma}^{-1}(\img_i)$. So 
$f_{\tau}(\img_i)=\img_{i-1}$ and $j\in 
f_{\tau}^{-1}(\img_i)$. It follows that $\img_i\in \supp(f_{\tau})$.
\medskip

If $\img_i$ fulfills case (\caseTwo), then by \cref{rem:case2},
$f_{\sigma}(\img_2)<\img_2$. 
There must therefore be some $j>\img_2$ such that $f_\sigma(j)=\img_2$,
since $\img_2 \in \supp(f_\sigma)$.
This implies that $\img_2\in \supp(f_{\tau})$. 
\medskip

If $\img_i$ fulfills case (\caseThree), then an argument similar  to that 
of case (\caseOne) on $\img_{i+1}$ shows that $\img_{i+1}\in \supp(f_{\tau})$.
\medskip


Finally, suppose $\img_i$ fulfills case (\caseFour).
\smallskip 

\textbf{Claim:} $f_{\sigma}(\img_{i+1})=\img_i$. \emph{Proof of claim:}
Assume the contrary. Then there is some $h<i$ such that 
$f_{\sigma}(\img_{i+1})=\img_h$, 
since  
\[
f_{\sigma}(\img_{i+1}) \in \{1,2,3,\ldots, \img_i,\img_{i+1}\},
\]
and  
$f_{\sigma}(\img_{i+1})<\img_{i+1}$.
%
Consider $\defin{t} \coloneqq i-1$. 
By our assumption, $i$ being the smallest with 
$\img_i$ satisfying case (\caseFour), no $\img_k$, for $2\leq 
k\leq t$, satisfies any one of the three 
conditions and for $k=i$, $\img_i$ does not satisfy 
both case (\caseOne) and case (\caseTwo).
By \cref{lem:fretherDIsWellDefined} either $f_{\sigma}$ starts as 
\begin{equation}\label{eq:imPreservingStep}
1\,1\,2\,3\,\dotsc\,\img_i{-}1, \quad \text{ where }\quad  \img_i{-}1= \img_{i-1}, 
\end{equation}
or 
\[
f_{\sigma}=1\,1\,2\,3\,\dotsc\,\img_{\ell-1}\,\img_\ell\,\dotsc\,\img_\ell
\]
with $f_{\sigma}^{-1}(\img_\ell)=\{\img_\ell+1,\img_\ell+2,\dotsc,n\}$. 
%
However, the second case is impossible since 
$\img_i<\img_\ell$ and $|f_\sigma^{-1}(\img_i)|\geq 2$,
and thus $f_{\sigma}$ is as in \eqref{eq:imPreservingStep}.

Since $f_\sigma(\img_i+1)=\img_i$ and  $h<i$, $f_{\sigma}$ starts as
\[
1\,1\,2\,3\,\dotsc\,\img_{i-1}\,\img_i
\]
and 
$f_{\sigma}(\img_h+1)=\img_h$. 
Then we have 
$f_{\sigma}(\img_h)<\img_h<\img_\ell$ and 
$\{\img_h+1\} \subsetneq f_{\sigma}^{-1}(\img_h)$.
%
This is a contradiction with $\img_i$ being the smallest 
with this property, cf.~\eqref{eq:case34Conditions}.
Therefore, $f_{\sigma}(\img_{i+1})=\img_i$, and the claim is proved.
\medskip 

We can now conclude $\img_i\in \supp(f_{\tau})$ in case (\caseFour). 
This proves the first equality in \eqref{eq:imPreserving}
and the second equality follows from \cref{prop:excAsIm}. 
\end{proof}

 
 
 
\begin{lemma}\label{lem:rlmPreserving}
For $f_{\tau} = \fretherD(f_{\sigma})$ we have 
\begin{equation}
 \RLMv(f_{\sigma}) = \RLMv(f_{\tau}) \text{ and }
 \RLMv(\sigma) = \RLMv(\tau).
\end{equation}
\end{lemma}
\begin{proof}
Let $f_\sigma\in \DSEF_n$ and $f_\tau=\fretherD(f_\sigma)$. 
If $f_\sigma$ is matchless, then $f_\tau=f_\sigma$ and  $\RLMv(f_\sigma)=\RLMv(f_\tau)$.
Suppose $f_\sigma$ is non-matchless, so that one of the four cases applies.

\begin{enumerate}[leftmargin=2cm]
 
	\item[\textbf{Case \caseOne:}] 
	There exists $i\geq 2$ such that $f_\sigma(\img_i)=\img_i$ and 
	$f_\tau(\img_i)=\img_{i-1}$. 
	Moreover, there is some $j>\img_i$ such that $f_\tau(j)=f_\sigma(j)=\img_i$. 
	The property of $\img_i$ being a right-left-minimum in $f_\sigma$
	as well as $f_\tau$ determined either at the position $j$
	or to the right of $j$. So, replacing $\img_i$ at position
	$\img_i$ in $f_\sigma$, preserves $\img_i$
	being a right-left-minima (or not).
	
	\begin{itemize}
		\item If $i\geq 3$, then $\img_i$ is the leftmost occurrence of $\img_{i-1}$ in $f_\tau$,
		since $i$ is the smallest such that $f_\sigma(\img_{i})=\img_i$.
		Since $\supp(f_\sigma)=\supp(f_\tau)$, there is some $k>\img_i$
		such that $f_\tau(k)=f_\sigma(k)=\img_{i-1}$.
		%
		So, $\RLMv(f_\sigma)=\RLMv(f_\tau)$.
		
		\item If $i=2$, then $\img_{i-1}=\img_1=1\in\RLMv(f_\tau)$. Since $1\in \RLMv(f_\sigma)$, then $\RLMv(f_\sigma)=\RLMv(f_\tau)$.
	\end{itemize}
	
	\item[\textbf{Case \caseTwo:}]
	By \cref{rem:case2}, $i=2$ and $f_\sigma(\img_2)=\img_1=1$
	and $f_\tau(\img_2)=\img_2$. 
	Moreover, there is some $k>\img_2$
	such that $f_\tau(k)=f_\sigma(k)=\img_2$.
	Since the right-left-minimum property of $\img_2$
	is determined at or to the right of the $k^\thsup$
	position and $1\in \RLMv(f_\tau)$, we have $\RLMv(f_\sigma)=\RLMv(f_\tau)$.
	
	
	\item[\textbf{Case \caseThree:}]
	There is some $i\geq 2$ such that $f_\sigma(\img_{i+1})=\img_{i+1}$
	and $f_\tau(\img_{i+1})=\img_{i}$.
	%
	We claim that $\img_i\in\RLMv(f_\sigma)$.
	Otherwise, there are $r<i$ and $s>j$,
	such that $f_\sigma(s)=\img_r$, where $j$ is the rightmost position of $\img_i$ in $f_\sigma$.
	%
	But now, $|f_\sigma^{-1}(\img_i)| \geq 2$
	so \eqref{eq:case34Conditions} is fulfilled.
	This contradicts the choice of $i$ being minimal and the claim follows.
	
	Since $\img_{i+1}$ being a right-left-minimum is determined at some other
	position $k>\img_{i+1}$ where $f_\tau(k)=f_\sigma(k)=\img_{i+1}$,
	we can conclude that $\RLMv(f_\sigma)=\RLMv(f_\tau)$.
	
	\item[\textbf{Case \caseFour:}]
	There is some $i\geq 2$ such that $f_\sigma(\img_{i+1})=\img_i<\img_{i+1}$
	and $f_\tau(\img_{i+1})=\img_{i+1}$.
	We also know that $\img_i\in\RLMv(f_\sigma)$. 
	Since $\img_{i+1}>\img_i$, we have $\img_i\in\RLMv(f_\tau)$.
	
	Now, $\img_{i+1}$ being a right-left-minimum
	is determined at some position $k>\img_{i+1}$ where 
	$f_\tau(k)=f_\sigma(k)=\img_{i+1}$.
	Hence, we can conclude that $\RLMv(f_\sigma)=\RLMv(f_\tau)$.
	\end{enumerate}
\end{proof}


\begin{lemma}\label{lem:signReversing}
If $f_{\sigma}\in \DSEF^*_n$, then 
\[
 \aexc(\fretherD(f_{\sigma})) \in \left\{ \aexc(f_{\sigma})-1, \aexc(f_{\sigma})+1 \right\}.
\]
Moreover, if $f_\tau = \fretherD(f_{\sigma})$,
then $\sigma$ and $\tau$ have different parity.
\end{lemma}
\begin{proof}
The first statement follows from the fact that $\fixMap_i$
and $\unfixMap_i$ decrease and increase, respectively, 
the number of strict anti-excedances by one.
The second statement follows from the first by \cref{prop:evenSEF}.
\end{proof}



\begin{lemma}\label{lem:involution}
	The map $\fretherD:\DSEF_n\longrightarrow \DSEF_n$
	is an involution.
\end{lemma}
\begin{proof}
	Let $f_\sigma\in \SEF_n$ and set
	$f_\tau=\fretherD(f_\sigma)$. For matchless $f_\sigma$, there is nothing to show,
	so suppose there is some $r \in \{2,3,\dotsc, \ell\}$ such that
	\[
	f_\tau=\fixMap_r(f_\sigma) \text{ or } f_\tau=\unfixMap_r(f_\sigma).
	\]
Now let $i$ be the smallest in $\{2,3,\dotsc, \ell\}$
fulfilling one of (\caseOne,\caseTwo,\caseThree,\caseFour) for $f_\sigma$.
	
\begin{enumerate}[leftmargin=1.6cm]
 	\item[\textbf{Case \caseOne:}] 
 	  We have $f_\sigma(\img_i)=\img_i$, 
 	  and $f_\tau=\unfixMap_i(f_\sigma)$.
		
		\begin{itemize}[leftmargin=-1cm]
			\item
			\emph{Suppose that $\vert f_\tau^{-1}(1)\vert \geq 3$}.
			We shall show that $f_\tau^{-1}(1) \neq f_\sigma^{-1}(1)$,
			by considering a few subcases.
			
			\medskip 
			\noindent 
			\textbf{Case $i=2$:}
			We  have $\img_2\in f_\tau^{-1}(1) \setminus f_\sigma^{-1}(1)$
			and the inequality is established.
			
		  There is now some $h>\img_2$ such that $f_\tau(h)=f_\sigma(h)=\img_2$,
		  so $f_\tau$ satisfies the conditions for case (\caseTwo) with $i'=2$. 
			It follows that $\fretherD(f_\tau)=f_\sigma$.
			
			\medskip 
			\noindent 
			\textbf{Case $i\geq 3$ and $|f_\sigma^{-1}(1)| = | f_\tau^{-1}(1)| \geq 3$:}
			This is impossible, since otherwise
			we would be in case (\caseTwo) for some $\img_j < \img_i$.
			
			\medskip 
			\noindent 
			\textbf{Case $i\geq 3$ and $|f_\sigma^{-1}(1)| \neq | f_\tau^{-1}(1)|$:} 
			%This is also impossible since $f_\sigma^{-1}(1) \neq  f_\tau^{-1}(1)$ implies $i=2$. 
			But then, there must be some $j\in [n]$ such that $f_\tau(j) = 1$ and $f_\sigma(j)\neq 1$
			(because $f_\sigma^{-1}(1) \subseteq  f_\tau^{-1}(1)$). However, $j=\img_i$ since $f_\tau$ and $f_\sigma$ differs
			only on position $\img_i$. So, $f_\tau(\img_i)=1=\img_{i-1}$  and this only happens if $i=2$,
			thus this case is empty.
			
			\medskip
			
			\item \emph{Suppose that $\vert f_\tau^{-1}(1)\vert = 2$}.
			%
			Now $f_\tau^{-1}(1)=f_\sigma^{-1}(1)$ because  
			we always have $f_\sigma^{-1}(1)\subseteq f_\tau^{-1}(1)$ in 
			(\caseOne) and (\caseThree),
			and since $|f_\sigma^{-1}(1)| \geq 2$, we must have equality.
			%
			Now, suppose $f_\sigma(2)=2$, then $i=2$ (because $f_\sigma \in \caseOne$)
			and $f_\tau = \unfixMap_2(f_\sigma)$, which would lead to
			$f_\tau^{-1}(1) \neq f_\sigma^{-1}(1)$.
			Hence, $f_\sigma(2)=1$ and 
			\[
			  f_\tau^{-1}(1)=f_\sigma^{-1}(1)=\{1,2\}.
			\]
			With some little work, one sees that this forces $i \geq 3$.
			
			Moving on, $f_\sigma(3) \in \{2,3\}$. We shall now treat each case separately.
			
			\medskip 
			\noindent 
			\textbf{Case $f_\sigma(3)=2$}. 
			This implies $ f_\tau(3)=2$ and $\img_2=2$.
			Now, $f_\sigma^{-1}(2)=\{3\}$, otherwise
			we have $f_\sigma$ belonging to (\caseFour) for $i=2$. 
			Now consider $f_\sigma(4)$; we may already conclude that
			\[
			 f_\sigma(4) \notin \{1,2\}, \text{ so } f_\sigma(4) \in \{3,4\}.
			\]
			Again, we shall consider the two cases separately, starting with $f_\sigma(4) = 4$.
			Here, we must have $\img_i=4$ for $i=3$ (if $3 \notin \supp(f_\sigma)$)
			or $i=4$ (if $3 \in \supp(f_\sigma)$).
			%
			It follows that $f_\sigma$ and $f_\tau$ are on one of the following two forms:
			\begin{equation}\label{eq:case1SomeSubcase}
			\begingroup % keep the change local
			\setlength\arraycolsep{2pt}
			\begin{cases}
			f_\sigma &= 
			\begin{pmatrix}
			1 & 2 & 3 & 4 & \dotsc & j & \dotsc \\
			1 & 1 & 2 & \mathbf{4} & \dotsc & 4 & \dotsc
			\end{pmatrix} \\
			f_\tau &= 
			\begin{pmatrix}
			1 & 2 & 3 & 4 & \dotsc & j & \dotsc \\
			1 & 1 & 2 & \mathbf{2} & \dotsc & 4 & \dotsc
			\end{pmatrix}
			\end{cases}
			\;
			\text{ or }
			\;
			\begin{cases}
			f_\sigma &= 
			\begin{pmatrix}
			1 & 2 & 3 & 4 & \dotsc & j & \dotsc & k & \dotsc \\
			1 & 1 & 2 & \mathbf{4} & \dotsc & 4 & \dotsc & 3 & \dotsc
			\end{pmatrix}\\
			f_\tau &= 
			\begin{pmatrix}
			1 & 2 & 3 & 4 & \dotsc & j & \dotsc & k & \dotsc \\
			1 & 1 & 2 & \mathbf{3} & \dotsc & 4 & \dotsc & 3 & \dotsc
			\end{pmatrix},
			\end{cases}
			\endgroup
			\end{equation}
			where $j$ is the index of the second $4$ in $f_\sigma$,
			and $k$ is the index\footnote{We either have $j>k$ or $j<k$; it will not matter for us.}
			of the first $3$ (if present).
			We now see that $f_\tau$ satisfies (\caseFour) for $i'=i-1$ in both cases of \eqref{eq:case1SomeSubcase},
			and thus $\fretherD(f_\tau)=f_\sigma$.
			
			\medskip 
			It remains to treat the case $f_\sigma(4)=3$.
			Now, we must have $f_\sigma^{-1}(3)=\{4\}$, otherwise
			$f_\sigma$ belongs to (\caseThree) for $i=3$. 
			%
			Repeating this argument inductively for $i=4,5,\dotsc,$
			we find that $f_\sigma$ must start as
			%
			\begingroup % keep the change local
			\setlength\arraycolsep{2pt}
			\begin{align*}
			f_\sigma = 
			\begin{pmatrix}
			1 & 2 & 3 & 4 & 5& \dotsc & \img_i-1\; & \img_i & \dotsc \\
			1 & 1 & 2 & 3 & 4& \dotsc & \img_i-2\; & \img_i & \dotsc   
			\end{pmatrix}.
			\end{align*}
			\endgroup
			There are now two options; either $m_i-1 \notin \supp(f_\sigma)$, or  $m_i-1 \in \supp(f_\sigma)$.
			This leads to the two corresponding situations shown in 
			\eqref{eq:case1SomeSubcase2a} and \eqref{eq:case1SomeSubcase2b}:
			%
			%
			\begingroup % keep the change local
			\begin{equation}\label{eq:case1SomeSubcase2a}
			\setlength\arraycolsep{2pt}
			\begin{cases}
				f_\sigma &= 
				\begin{pmatrix}
				1 & 2 & 3 & 4 & \dotsc & \img_i-1\;& \img_i & \dotsc \\
				1 & 1 & 2 & 3 & \dotsc & \img_i-2\;& \img_i & \dotsc 
				\end{pmatrix} \\
				f_\tau &= 
				\begin{pmatrix}
					1 & 2 & 3 & 4 & \dotsc & \img_i-1\;& \img_i   & \dotsc \\
					1 & 1 & 2 & 3 & \dotsc & \img_i-2\;& \img_i-2 & \dotsc 
				\end{pmatrix}
			\end{cases}
			\end{equation}
			where $m_{i-1} = m_i-2$, or
			%
			\begin{equation}\label{eq:case1SomeSubcase2b}
			\begin{cases}
			f_\sigma &= 
			\begin{pmatrix}
			1 & 2 & 3 & 4 & \dotsc & \img_i-1\;& \img_i & \dotsc & k     & \dotsc\\
			1 & 1 & 2 & 3 & \dotsc & \img_i-2\;& \img_i & \dotsc & m_i-1 & \dotsc
			\end{pmatrix} \\
			f_\tau &= 
			\begin{pmatrix}
			1 & 2 & 3 & 4 & \dotsc & \img_i-1\;& \img_i &   \dotsc & k     & \dotsc\\
			1 & 1 & 2 & 3 & \dotsc & \img_i-2\;& \img_i-1 & \dotsc & m_i-1 & \dotsc
			\end{pmatrix},
			\end{cases}
			\end{equation}
			\endgroup
			where $m_{i-1} = m_i-1$. In both \eqref{eq:case1SomeSubcase2a} and \eqref{eq:case1SomeSubcase2b},
			$f_\tau$ belongs to (\caseFour) for $i'=i-1$, and $\fretherD(f_\tau)=f_\sigma$.
			
			\medskip 
			\noindent 
			\textbf{Case $f_\sigma(3)=3$}. 
			Here, we have
			\[
				i=3, \text{ and } \img_3=3 \text{ (since $i\geq 3$)}.
			\]
			Then there must be some $j>3$ such that $f_\sigma(j)=2$,
			and it follows that $f_\tau(3)=\img_2=2$. 
			Moreover, there must be some $k \geq 4$, with $f_\sigma(k)=3$. 
			Hence, we are in the situation
			\begin{equation}\label{eq:case1SomeSubcase3}
			\setlength\arraycolsep{2pt}
			\begin{cases}
				f_\sigma &= 
				\begin{pmatrix}
				1 & 2 & 3 &  \dotsc & j & \dotsc  & k & \dotsc\\
				1 & 1 & 3 &  \dotsc & 2 & \dotsc  & 3 & \dotsc
				\end{pmatrix} \\
				f_\tau &= 
				\begin{pmatrix}
				1 & 2 & 3 &  \dotsc & j & \dotsc  & k & \dotsc\\
				1 & 1 & \mathbf{2} &  \dotsc & 2 & \dotsc  & 3 & \dotsc
				\end{pmatrix}
			\end{cases}
			\end{equation}
			where $f_\tau$ belongs to case (\caseFour) for $i'=2$.
			Therefore, $\fretherD(f_\tau)=f_\sigma$.
		\end{itemize}
		
		\item[\textbf{Case \caseTwo:}] 
		By \cref{rem:case2}, $i=2$, $f_\tau = \fixMap_2(f_\sigma)$.
		We have that $f_\tau(\img_2)=\img_2$ and $f_\tau$ belongs to (\caseOne) for $i'=i$,
		so $\fretherD(f_\tau)=f_\sigma$.
		
		
		
	\item[\textbf{Case \caseThree:}] In this case,
	  we have \eqref{eq:case34Conditions}, stating that 
		\[
		f_{\sigma}(\img_i) < \img_i < \img_\ell,\quad
		f^{-1}_\sigma(1) = \{1,2\}, \text{ and }
		\{\img_i + 1\}\subsetneq f^{-1}_\sigma(\img_i),
		\]
	  and that $f_\sigma(\img_{i+1})=\img_{i+1}$.
		We must have that $f_\sigma = 112 \dotsc$, otherwise $f_\sigma$ belongs to
		(\caseOne) for $\img_i=3$. Recall also that $f_\tau = \unfixMap_{i+1}(f_\sigma)$.
		Our goal is to show that $f_\tau$ belongs to $(\caseFour)$ for $i'=i$.
		
		Firstly, $f_\tau(\img_i)<\img_i<\img_\ell$
		and we can easily show that $f_\tau^{-1}(1)=\{1,2\}$.
		
		%
		Since 
		$f_\sigma(\img_i+1)=\img_i$ and  
		$f_\sigma(\img_{i+1})=\img_{i+1}$, we have $\img_i+1\neq \img_{i+1}$. 
		This implies that $f_\tau(\img_i+1)=f_\sigma(\img_i+1)=\img_i$.
		
		Finally, we also have that $|f_\tau^{-1}(\img_i)| \geq 3$,
		since $|f_\sigma^{-1}(\img_i)|\geq 2$ and $f_\tau(\img_{i+1})=\img_i$. 
		Hence, $f_\tau$ belongs to (\caseFour) for $i'=i$,
		and it follows that $\fretherD(f_\tau)=f_\sigma$.
		
		\item[\textbf{Case \caseFour:}] 
		Again, we have \eqref{eq:case34Conditions} and $f_\sigma(\img_{i+1})<\img_{i+1}$.
		Moreover, in the proof of \cref{lem:imPreserving},
		we proved that $f_\sigma(\img_{i+1})=\img_{i}$.
		Recall, $f_\tau = \fixMap_{i+1}(f_\sigma)$.
		We shall now consider two subcases.
		%
		\begin{itemize}
			\item \emph{Suppose $\img_i+1=\img_{i+1}$.} 
			Then $f_\tau(\img_i)=f_\sigma(\img_i)<\img_i$, but 
			$f_\tau(\img_i+1)=f_\tau(\img_{i+1})=\img_{i+1}\neq 
			\img_i$. So $f_\tau$ belongs to (\caseOne) for $i'=i+1$,
			and $\fretherD(f_\tau)=f_\sigma$.
			
			\item \emph{Suppose now that $\img_i+1< \img_{i+1}$.}
			Then $f_\tau(\img_i)=f_\sigma(\img_i)<\img_i<\img_\ell$
			and $f_\sigma^{-1}(1)=\{1,2\}$. % since $i\geq 2$ implies $i+1\geq 3$.
			Moreover, $\img_i+1 \neq  \img_{i+1}$, we have that 
			$f_\tau(\img_i+1)=f_\sigma(\img_i+1)=\img_i$,
			and we also have $f_\tau(\img_{i+1})=\img_{i+1}$.
			
			If $|f_\sigma^{-1}(\img_i)|=2$, 
			then $f_\tau$ lies in case (\caseOne) for $i'=i+1$.
			
			If $\vert f_\sigma^{-1}(\img_i)\vert\geq 3$, 
			then $f_\tau$ belongs to case (\caseThree) for $i'=i$.
			
			In both cases, $\fretherD(f_\tau)=f_\sigma$.
		\end{itemize}
	\end{enumerate}
\end{proof}


\begin{remark}
In \cref{tab:casesCombinations}, we give an overview under what circumstances 
a subexcedant function belonging to a case, is mapped to a different case.
When $f_\tau = \fretherD(f_\sigma)$ we have the following combinations, under the 
conditions described in the cells of the table.
%
\begin{table}[!ht]
\begin{center}
\begin{tabular}{lp{2cm}p{2cm}p{2cm}p{2cm}}
 \; & $f_\sigma\in\caseOne$ & $f_\sigma\in\caseTwo$ & $f_\sigma\in\caseThree$ & $f_\sigma\in\caseFour$ \\
   \toprule 
$f_\tau\in\caseOne$   & $\emptyset$  & Always & $\emptyset$ & See proof \\
$f_\tau\in\caseTwo$   & $|f_\tau^{-1}(1)|\geq 3$  & $\emptyset$ & $\emptyset$ & $\emptyset$\\
$f_\tau\in\caseThree$ & $\emptyset$  & $\emptyset$ & $\emptyset$ & See proof\\
$f_\tau\in\caseFour$  & $|f_\tau^{-1}(1)| =2$  & $\emptyset$ & Always & $\emptyset$
\end{tabular}
\caption{
The conditions on $\sigma$ and $\tau$ for case $(\caseFour)$
are more involved, and we refer to \cref{lem:involution} for the exact conditions.
}\label{tab:casesCombinations}
\end{center}
\end{table}
\end{remark}


We finish this subsection with all properties proved about $\fretherD$.
\begin{corollary}\label{cor:fretherDProps}
The map $\fretherD:\DSEF_n\longrightarrow \DSEF_n$
is an involution with the following properties.
%
\begin{enumerate}[label=(\roman*)]
\item 
	The image is preserved, $\supp(\fretherD(f_\sigma )) = \supp( f_\sigma )$.
\item 
	If $f_\tau = \fretherD(f_{\sigma})$, then 
	$
		\EXCv(\tau) = \EXCv(\sigma).
	$
\item 
	The set of right-to-left minima is preserved,
	$\RLMv( \fretherD(f_\sigma ) ) = \RLMv( f_\sigma )$.
\item 
	Whenever $f_\sigma \in \DSEF^*_n$, 
	\[
	\sgn(\fretherD(f_\sigma )) = -\sgn( f_\sigma ).
	\]
\end{enumerate} 
\end{corollary}
We may now consider $\fretherD$ as an involution on derangement, \defin{$\fretherD:\dArr_n \to \dArr_n$},
by setting
\[
  \defin{\fretherD(\sigma)} \coloneqq (\sefToPerm \circ \fretherD \circ \sefToPerm^{-1})(\sigma)
  \text{ for $\sigma \in \dArr_n$.}
\]

\subsection{Consequences}

Before stating the main theorem, we shall first 
introduce two auxiliary involution on $ \symS_n$,
and prove some of its properties.
Let $\defin{\flipMap} : \symS_n \to \symS_n$ be the map 
\[
\defin{\flipMap(\sigma)(k)} \coloneqq n+1-\sigma(k) \quad \text{ for } \quad k\in [n],
\]
and let $\defin{\zetaMap} : \symS_n \to \symS_n$ be the composition
$
 \defin{\zetaMap} \coloneqq \flipMap^{-1} \circ (\;\cdot\;)^{-1} \circ \flipMap.
$
In other words,
\[
  \defin{\zetaMap(\sigma)(k)} \coloneqq n+1-\sigma^{-1}(n+1-k) \quad \text{ for } \quad k\in [n].
\] 
\begin{lemma}\label{lem:conjugateProps}
The map $\zetaMap$ is an involution, and
\begin{align*}
 \EXCv(\pi) &= \{ n+1-k : k \in \EXCi( \zetaMap(\pi)) \}, \\
 \RLMv(\pi) &= \{ n+1-k : k \in \RLMi( \zetaMap(\pi)) \},\\
 \FIX(\pi)  &= \{ n+1-k : k \in \FIX( \zetaMap(\pi)) \},\\
 \inv(\pi)  &= \inv(\zetaMap(\pi)).
\end{align*}
In particular, $\zetaMap$ restricts to a sign-reversing
involution $\zetaMap:\dArr_n \to  \dArr_n$.
\end{lemma}
\begin{proof}
It follows immediately from the definition that $\zetaMap$ is an involution. 

For the first property, let $i\in [n]$ and set $j \coloneqq n+1-i$.
We then see that
\begin{align*}
 i\in\EXCv(\pi) &\iff \pi^{-1}(i)<i \\
                &\iff n+1-\pi^{-1}(i)>n+1-i \\
                &\iff n+1-\pi^{-1}(n+1-j)>j \\
                &\iff \zetaMap(\pi)(j)>j \\
                &\iff i\in \{n+1-k:k \in \EXCv(\zetaMap(\pi)) \}.
\end{align*}
Now for the right-to-left minima. Again with $j \coloneqq n+1-i$, we have 
\begin{align*}
i\in \{n+1-k:k \in \RLMi(\zetaMap(\pi)) \} &\iff 
\zetaMap(\pi)(j)<\zetaMap(\pi)(t) \text{ whenever } j<t \leq n \\
&\iff \pi^{-1}(n+1-j)>\pi^{-1}(n+1-t) \text{ whenever } j<t \leq n \\
&\iff \pi^{-1}(i)>\pi^{-1}(k) \text{ whenever } k\in[i-1] \\
&\iff \text{every $k\in[i-1]$ lies to the left of $i$ in $\pi$}\\
&\iff i\in\RLMv(\pi).
\end{align*}
%
Similarly, with $i \in [n]$, $j \coloneqq n+1-i$,
\begin{align*}
i\in\FIX(\pi) &\iff \pi^{-1}(i)=i \\
&\iff \pi^{-1}(n+1-j)=n+1-j\\
&\iff n+1-\pi^{-1}(n+1-j)=j\\
&\iff \zetaMap(\pi)(j)=j\\
&\iff j\in\FIX(\zetaMap(\pi)).
\end{align*} 
This last property also shows that $\zetaMap$ is an involution on $\dArr_n$.

Finally, we have that
\begin{align*}
	\inv(\zetaMap(\pi)) &= |\{ (i,j) : 1\leq i < j \leq n \text{ such that } \zetaMap(\pi)(i)>\zetaMap(\pi)(j) \}|\\
	&=|\{ (i,j) : 1\leq i < j \leq n \text{ such that }\pi^{-1}(n+1-i)<\pi^{-1}(n+1-j) \}|\\
	&=|\{ (n+1-k,n+1-l) : 1\leq l < k \leq n \text{ such that }\pi^{-1}(k)<\pi^{-1}(l) \}|\\
	&=|\{ (l',k') : 1\leq l' < k' \leq n \text{ such that }\pi^{-1}(k')<\pi^{-1}(l') \}|\\
	&=\inv(\pi),
\end{align*}
so $\zetaMap$ preserves the number of inversions.
In particular, $\zetaMap$ preserves the sign.
\end{proof}


We are now ready to prove the main theorem in this paper.
\begin{theorem}\label{thm:theMainTheorem}
 We have that 
 \begin{equation}\label{eq:mainResultInText}
 \sum_{\pi \in \dArr_n} (-1)^{\inv(\pi)} 
 \xvec_{\RLMv(\pi)} \yvec_{ \EXCv(\pi)} 
 =
 (-1)^{n-1} \sum_{j=1}^{n-1} x_1 \dotsm x_j \cdot y_{j+1}\dotsm y_{n}.
 \end{equation} 
Moreover, 
 \begin{equation}\label{eq:mainResultInText2}
 \sum_{\pi \in \dArr_n} (-1)^{\inv(\pi)} 
 \xvec_{\RLMi(\pi)} 
 \yvec_{\EXCi(\pi)}
 =
 (-1)^{n-1} \sum_{j=1}^{n-1} y_1 \dotsm y_j \cdot x_{j+1}\dotsm x_{n}.
 \end{equation} 
\end{theorem}
\begin{proof}
By \cref{prop:excAsIm,prop:rlmForSEF,prop:evenSEF}, we have that
\begin{equation}\label{eq:mainThmSumOverSEF}
 \sum_{\pi \in \dArr_n} (-1)^{\inv(\pi)} 
 \xvec_{\RLMv(\pi)} \yvec_{ \EXCv(\pi)} 
 =
\sum_{f_\sigma \in \DSEF_n} (-1)^{\aexc(f_\sigma)} 
 \xvec_{\RLMv(f_\sigma)} \yvec_{[n] \setminus  \supp(f_\sigma)}. 
\end{equation}
By applying the involution $\frether$ and all 
using the properties listed in \cref{cor:fretherDProps},
all terms in the right-hand side with $f_\sigma \in \DSEF^*_n$
cancel, so \eqref{eq:mainThmSumOverSEF} equals
\begin{equation}
\sum_{\substack{ f_\sigma \in  \DSEF_n \\ \text{$f_\sigma$ matchless}}}
(-1)^{\aexc(f_\sigma)} \xvec_{\RLMv(f_\sigma)} \yvec_{[n] \setminus  \supp(f_\sigma)}. 
\end{equation}
Finally, by \cref{lem:propertiesOfMatchless}, this sum is equal to
\begin{equation}
\sum_{k=1}^{n-1}
(-1)^{n-1} \xvec_{[k]} \yvec_{[n] \setminus [k]},
\end{equation}
which is the right-hand side of \eqref{eq:mainResultInText}.
\medskip

Let $\defin{\rho(S)}\coloneqq \{n+1-s : s\in S\}$ whenever $S\subseteq [n]$.
By applying the change of variables
$x_i \mapsto x_{n+1-i}$, $y_j \mapsto y_{n+1-j}$ on both sides of 
\eqref{eq:mainResultInText}, we get
\begin{align*}
 \sum_{\pi \in \dArr_n} (-1)^{\inv(\pi)} 
 \xvec_{\rho(\RLMv(\pi))}
  \yvec_{\rho(\EXCv(\pi))}
&= (-1)^{n-1} \sum_{j=1}^{n-1} x_n \dotsm x_{n-j} \cdot y_{n-j}\dotsm y_{1} \\
&= (-1)^{n-1} \sum_{j'=1}^{n-1} x_{j'+1} \dotsm x_{n} \cdot y_{1}\dotsm y_{j'}.
\end{align*}
%
Now by \cref{lem:conjugateProps},
\begin{align*}
\sum_{\pi \in \dArr_n} (-1)^{\inv(\pi)} 
\xvec_{\rho(\RLMv(\pi))}
\yvec_{\rho(\EXCv(\pi))}
\!
=
\!
\sum_{\pi \in \dArr_n} (-1)^{\inv(\zetaMap(\pi))} 
\xvec_{\RLMi(\zetaMap(\pi))}
\yvec_{\EXCi(\zetaMap(\pi))}.
\end{align*}
Since $\zetaMap$ sends $\dArr_n$ to $\dArr_n$, 
the last sum must be exactly the right-hand side of \eqref{eq:mainResultInText2}, and we are done.
\end{proof}



\begin{corollary}
By letting $x_j \to 1$, $z_j \to 1$ and $y_j \to t$, we have that 
\begin{equation}
 \sum_{\pi \in \dArr_n} (-1)^{\inv(\pi)} t^{\exc(\pi)}
 =
 (-1)^{n-1}(t+t^2+\dotsb+t^{n-1}).
 \end{equation}
By comparing coefficients of $t^k$,
we get \eqref{eq:mainEquation}, the main result of \cite{MantaciRakotondrajao2003}, 
%
In a similar manner,
\begin{equation}
\sum_{\pi \in \dArr_n} (-1)^{\inv(\pi)} t^{\rlm(\pi)}
 =
 (-1)^{n-1}(t+t^2+\dotsb+t^{n-1}).
\end{equation}
\end{corollary}


\section{A proof using generating functions}\label{sec:genFunc}


We shall first define an involution
$\iota:\symS_n \to \symS_n$ such that $\iota$
\begin{enumerate}
 \item preserves the set of excedances, $E$,
 \item $\iota$ changes sign of non-fixed elements,
 \item all fixed-points of the involution has sign $(-1)^{|E|}$,
 \item there is a \emph{unique} fixed-point with excedance set $E \subseteq [n-1]$.
\end{enumerate}
%
We shall now describe an such involution $\iota$,
which is essentially the one given in \cite{Mantaci1993}.


\begin{definition}
Let $E\subseteq [n-1]$ be the set of excedances in a permutation $\pi$ over $[n]$. 
Define a mapping $\defin{\iota}: \symS_n \rightarrow \symS_n$ by 
$\iota(\pi)=\pi'$,
where $\pi'$ is obtained from $\pi$ by swapping $\pi(l)$ and $\pi(m)$,
for $(l,m)=\max\{(i,j): 2\leq i<j\leq n$ and either $i, j\in E$ or $i,j\notin E\}$ 
with respect to lexicographical order, so that the set of excedances in $\pi'$ is $E$. 
It can be defined as $\iota(\pi)=\pi (l,m)$ if $\pi$ is in cycle form.
If there is no such $(l,m)$, then $\iota(\pi)=\pi$ and we say that $\pi$ is \defin{critical}.
\end{definition}
From the definition, it is clear that (1) and (2) hold.
We must show that (3) and (4) hold as well.


\begin{lemma}
Suppose $\pi \in \symS_n$, and that $i<j$ are excedances of $\pi$
where $\pi(i)=a>i$ and $\pi(j)=b>j$ for some $a>b$.
Then $\pi$ is not critical.

Similarly, if $i'<j'$ are anti-excedances of $\pi$
with $\pi(i')=a'\leq i'$ and $\pi(j')=b' \leq j'$ for some $a'>b'$,
then $\pi$ is not critical.
\end{lemma}
\begin{proof}
Note that swapping the excedances $a$ and $b$,
produces a permutation $\pi'$ with $\pi'(i)=b$ and $\pi'(j)=a$.
However, both $i$ and $j$ are excedance for $\pi'$, since $\pi'(i)=b>j>i$,
and $\pi'(j)=a>b>j$. Hence, $\pi$ is not critical as there is at least
one pair of entries where  we can perform a swap as in the definition of $\iota$.
A similar argument proves the second statement.
\end{proof}

The previous lemma implies that if $\{j_1,\dotsc,j_k\}$
are the excedances of a critical permutation $\pi$,
and $\{i_1,\dotsc,i_{n-k}\}$ are the anti-excedances,
then
\begin{equation}\label{eq:excSorted}
\begin{aligned}
 &j_1 < \pi(j_1) \leq j_2 < \pi(j_2) \leq j_3 < \pi(j_3) \leq \dotsb \leq j_k< \pi(j_k), \\
 &\pi(i_1) \leq i_1 < \pi(i_2) \leq i_2 < \pi(i_3) \leq i_3 <\dotsb < \pi(i_{n-k}) \leq i_{n-k}.
\end{aligned}
\end{equation}

Moreover, if $i_{n-k-1}<j_1$ for $\pi \in \symS_n$,
so that 
\[
  i_1 < i_2< \dotsb < i_{n-k-1}<j_1<j_2 < \dotsb < j_k <i_{n-k}
\]
then it follows directly that 
\begin{equation}\label{eq:specialCase}
  \pi = (n-k\quad n-k+1 \quad \dotsc \quad n-1\quad  n),
\end{equation}
and we note that $\inv(\pi)=|E|$.



The following is similar to an argument in \cite{Mantaci1993},
where a slightly different\footnote{We note that the original proof has a few typos.}
approach is taken.
\begin{proposition}\label{prop:mantaciInvolution}
Let $E=\{j_1,j_2,\dotsc,j_k\} \subseteq [n-1]$, and define the permutation $\pi_E \in \symS_n$
with excedance set $E$ via
\begin{align*}
    \pi_E(j_s) &=j_s+1 \quad &&\text{ for each } j_s\in E, \\
    \pi_E(i_r) &=i_{r-1}+1\quad &&\text{ for each } i_r\in [n]\setminus E, \,\,(j_0 \coloneqq 0).
\end{align*}
Then $\pi_E$ is the unique critical permutation in $\symS_n$ with excedance set $E$,
and $\inv(\pi_E)=|E|$.
\end{proposition}
\begin{proof}
We first show that $\pi_E$ is critical, so assume it is not. 
Then there is some $(l,m)$ such that swapping $\pi(l)$ and $\pi(m)$,
produces a $\pi'$ from $\pi$ with the same set of (anti)excedances.
But swapping $j_s+1$ and $j_{s'}+1$ for some $s < s'$, will result in 
$\pi'(j_{s'}) = j_s+1 < j_{s'}+1 \leq j_{s'}$,
so the set of excedances is not preserved.
A similar argument shows that we cannot swap a pair of anti-excedances
either.

Now that we have established that $\pi_E$ is critical, we must show that
there are no other critical permutations in $\symS_n$ with $E$ as excedance set.

We proceed by (strong) induction over $n$. The base case $n=1$ is trivial.

First we handle the case $|E|=n-1$ where $E = \{1,2,\dotsc,n-1\}$.
There is only one permutation with $E$ as excedance set,
namely $\pi = (1\;2\; \dotsc \; n)$ in cycle form, and this is exactly $\pi_E$,
where $\inv(\pi_E)=|E|$.

Suppose now that $|E|=k<n-1$, $E = \{j_1,\dotsc,j_k\}$, and $\{i_1,\dotsc,i_{n-k}\}$ be 
the set of anti-excedances of some critical permutation $\pi$.
Note that $i_{n-k}=n$.
Let $h \in [n]$ be the largest integer such that $j_h < i_{n-k-1}$.
If there is no such $h$, then $E = \{n-k,\dotsc,n-2,n-1\}$
and $\pi$ is of the form given in \eqref{eq:specialCase}
and the statement holds.
%
So now, there is some $m \in [n-k]$ such that $j_h+1=i_{n-k-m}$,
and the permutation $\pi$ has the following structure:
\[
 \pi =
 \begingroup
 \setcounter{MaxMatrixCols}{20}
 \setlength\arraycolsep{2pt}
 \scriptstyle{
 \begin{pmatrix}
 1      & 2      & \dotsb & j_h & i_{n-k-m}  & \dotsb   & i_{n-k-1}      & j_{h+1} & \dotsb    & j_{k} &  n \\
 \pi(1) & \pi(2) & \dotsb & \pi(j_h) & \pi(i_{n-k-m}) & \dotsb & \pi(i_{n-k-1}) & \pi(j_{h+1}) & \dotsb & \pi(j_{k})  & \pi(n)
\end{pmatrix}}.
\endgroup
\]
We have that
\[
j_{h+1} < \pi(j_{h+1}), \quad j_{h+2}<\pi(j_{h+2}), \quad j_{k}<\pi(j_{k}),
\]
and taking \eqref{eq:excSorted} into account, this is only possible if 
\[
 \pi(j_{h+1}) = j_{h+2},\quad \pi(j_{h+2}) = j_{h+3},\quad \dotsc, \quad \pi(j_{k}) = n.
\]
%
\textbf{Suppose now $\pi(j_h)>i_{n-k-1}$}. Then by \eqref{eq:excSorted},
either $\pi(j_h)=j_{h+1}$, or $h=k$ and $\pi(j_h)=n$.
In either case, we must have that
$\pi(n)=i_{n-k-1}$ and $\pi(i_{n-k-1}) < i_{n-k-1}$.
But this is not possible, as $\pi$ would not be critical (we could swap
the values at positions $n-k-1$ and $n$).
Hence, $\pi(j_h) \leq i_{n-k-1}$.
\medskip 

It follows that $\pi$ sends $\{1,2,\dotsc,i_{n-k-1}\}$ to itself
and that $\pi(n) = i_{n-k-1} + 1$. Hence,
the value of $\pi(s)$ is uniquely determined for all
$s > i_{n-k-1}$.
So, $\pi$ restricts to a critical permutation $\pi'$
acting on $[i_{n-k-1}]$.
By induction, $\pi'$ is uniquely determined by $E \cap [i_{n-k-1}]$ with 
so it follows that $\pi$ is unique and of the form $\pi_E$.
Also by induction, $\inv(\pi') = h = |E \cap [i_{n-k-1}]|$,
and finally 
\[
  \inv(\pi) = \inv(\pi') + (k-h) = h + (k-h) = |E|.
\]
\end{proof}




\begin{example}
 Let $n=4$ and consider all permutations with $2$ excedances.
 We have $7$ even permutations with two exedances,
 and $4$ odd permutations. The sign-reversing involution should therefore have 
 $\binom{3}{2}=3$ fixed-points, all with even sign $(-1)^2$.
 \begin{center}
\begin{tabular}{cc}
Even & Odd \\
\toprule
1342 & $\mathbf{1342}$\\
2143 & $\mathbf{2143}$\\
2314 & $\mathbf{2314}$\\
2431 & 2413  \\
3241 & 3142 \\
3412 & 3421 \\
4321 & 4312 \\
\end{tabular}
\end{center}
\end{example}


\cref{prop:mantaciInvolution} now immediately gives a bijective proof of the following result. 
The same was of course done in \cite{Mantaci1993} (albeit stated in a slightly different manner).
\begin{proposition}\label{lem:excEvenOdd}
Let $n \geq 1$, then
\begin{equation}\label{eq:excEvenOdd}
 \sum_{\substack{\pi \in \symS_n}} (-1)^{\inv(\pi)} \xvec_{\EXCi(\pi)} = 
 \prod_{j \in [n-1]} (1-x_j) = \sum_{E \subseteq [n-1]} (-1)^{|E|}\xvec_{E}.
\end{equation}
In particular, by setting all $x_i$ equal to $t$, 
we have
\[
 \sum_{\substack{\pi \in \symS^e_n}} t^{\exc(\pi)} 
 -
 \sum_{\substack{\pi \in \symS^o_n }} t^{\exc(\pi)} = (1-t)^{n-1}.
\]
\end{proposition}





\begin{proposition}\label{prop:main}
Let $n\geq 1$ and let $T \subseteq [n]$. 
Let $m \leq n$ be the largest integer not in $T$ and set
$
  E = \{1,2,\dotsc,m-1\} \setminus T.
$
Then
 \begin{equation}\label{eq:main}
 \sum_{\substack{\pi \in \symS_n  \\ T \subseteq \FIX(\pi)  }} 
 (-1)^{\inv(\pi)} \xvec_{\EXCi(\pi)} =  
 \prod_{j \in E} (1-x_j),
\end{equation}
where the empty product has value $1$.

Setting all $x_i$ to be $t$, we have
\begin{equation}
 \sum_{\substack{\pi \in \symS^e_n  \\ T \subseteq \FIX(\pi)  }} t^{\exc(\pi)}
 -
 \sum_{\substack{\pi \in \symS^o_n \\ T \subseteq \FIX(\pi)   }} t^{\exc(\pi)}
 = 
 \begin{cases}
 1 & \text{if $|T|=n$} \\
 (1-t)^{n-1-|T|} & \text{otherwise.}
 \end{cases}
\end{equation}
\end{proposition}
\begin{proof}
First note that $E = \emptyset$ if $T=[n]$ and \eqref{eq:main} is easy to verify,
so from now on, we may assume $|T| < n$.

By definition of $m$, we have that $T = T_1 \cup T_2$
where $T_1 \subseteq \{1,2,\dotsc,m-1\}$, and $T_2 = \{m+1,m+2,\dotsc,n\}$.
Hence, $|E|+|T_1| = m-1$ and $|T_2| = n-m$, and 
\[
|E| = n-1-|T_1|-|T_2| = n-1-|T|.
\]
Now suppose $\pi \in \symS_n$ is a permutation such that $T \subseteq \FIX(\pi)$.
We then construct $\pi' \in \symS_{n-|T|}$, by only considering the positions not in $T$,
and the relative ordering of the entries at these positions.
For example, if $n=9$, $T=\{2,4,6,8,9\}$,
$[n] \setminus T = \{1,3,5,7\}$ and
\[
  \pi = \underline{1} 2 \underline{7} 4 \underline{3} 6 \underline{5} 8 9,
  \text{ we have } \pi' = 1423
\]
since the relative ordering of $1,3,5$ and $7$ in $\pi$ is $1423$.
Observe that $\exc(\pi) = \exc(\pi')$ and $(-1)^{\inv(\pi)} = (-1)^{\inv(\pi')}$.

Hence, the sum in the left-hand side of \eqref{eq:main},
can be taken as a sum over permutations $\pi' \in \symS_{n-|T|}$,
but with a reindexing of the variables using values in $[n] \setminus T$.
Now, this sum can be computed using \cref{lem:excEvenOdd}
which finally gives \eqref{eq:main}.
Note that $m$ is the largest member of $[n] \setminus T$, 
so we do not get any variable with this index --- this corresponds
to the fact that the right-hand side of \eqref{eq:excEvenOdd}
only uses elements in $[n-1]$.
\end{proof}



\begin{theorem}\label{thm:mainGenFunc}
Let $n\geq 1$. Then
\begin{equation}\label{eq:derangementSum}
 \sum_{\pi \in \dArr_n } (-1)^{\inv(\pi)} \xvec_{\EXCi(\pi)} = 
 (-1)^{n-1} \sum_{j=1}^{n-1} x_{1}x_{2}\dotsm x_{j}.
\end{equation}
\end{theorem}
\begin{proof}
By inclusion-exclusion, we have the two identities
\begin{align*}
 \sum_{\substack{\pi \in \symS^e_n \\ \FIX(\pi)=\emptyset}} \xvec_{\EXCi(\pi)}
 &=
 \sum_{\substack{T \subseteq [n]  }}
 (-1)^{|T|} \sum_{\substack{\pi \in  \symS^e_n  \\ T \subseteq \FIX(\pi)   }} \xvec_{\EXCi(\pi)}\\
%
 \sum_{\substack{\pi \in \symS^o_n \\ \FIX(\pi)=\emptyset}} \xvec_{\EXCi(\pi)}
 &=
 \sum_{\substack{T \subseteq [n]   }}
 (-1)^{|T|} \sum_{\substack{\pi \in  \symS^o_n  \\ T \subseteq \FIX(\pi)   }} \xvec_{\EXCi(\pi)}.
\end{align*}
%
%
By taking the difference of these two identities, we get
\begin{align*}
 \sum_{\substack{\pi \in \dArr_n }} (-1)^{\inv(\pi)} \xvec_{\EXCi(\pi)}
 &=
\sum_{\substack{T \subseteq [n] }}
 (-1)^{|T|}
\left( 
 \sum_{\substack{\pi \in  \symS^e_n  \\ T \subseteq \FIX(\pi)   }} \xvec_{\EXCi(\pi)} 
 -
 \sum_{\substack{\pi \in  \symS^o_n  \\ T \subseteq \FIX(\pi)   }} \xvec_{\EXCi(\pi)}
\right). 
\end{align*}
By \cref{prop:main}, the difference in the right-hand side is equal to 
\[
  \prod_{ j \in [m_T-1]\setminus T } (1-x_j)
\]
where $m_T\leq n$ is the largest integer not in $T$.
We divide the terms depending on the value of $m_T$.
If $m_T=0$ then $T=[n]$ and the product is empty, so its value is $1$.
In general, the left hand side of \eqref{eq:derangementSum} is equal to
\[
(-1)^{n} + 
 \sum_{k=1}^n \sum_{\substack{T \subseteq [n] \\ m_T = k}} (-1)^{|T|}
 \prod_{ j \in [k-1]\setminus T } (1-x_j). 
 \]
By using $E = [k-1]\setminus T$, this can then be simplified to
 \[
 (-1)^{n} +
 \sum_{k=1}^n 
 \sum_{\substack{E \subseteq [k-1]}} (-1)^{n-1-|E|}
 \prod_{ j \in E } (1-x_j).
\]
Canceling the $k=1$ case with $(-1)^{n}$, and then shifting the index, we get
 \[
(-1)^{n-1}
 \sum_{k=1}^{n-1} 
 \sum_{\substack{E \subseteq [k]}} 
 \prod_{ j \in E } (x_j-1).
\]
Now,
\begin{align*}
(-1)^{n-1}
\sum_{k=1}^{n-1} 
\sum_{\substack{E \subseteq [k]}} 
\prod_{ j \in E } (x_j-1)
&=
(-1)^{n-1}
\sum_{k=1}^{n-1} 
\sum_{\substack{E \subseteq [k]}} 
\sum_{F \subseteq E} (-1)^{|E|-|F|}\xvec_F \\
&=
(-1)^{n-1}
\sum_{k=1}^{n-1} 
\sum_{\substack{F \subseteq [k]}} 
(-1)^{|F|}\xvec_F
\sum_{F \subseteq E \subseteq [k]} (-1)^{|E|}.
\end{align*}
The last sum vanish unless $F=[k]$, and we thus get that 
\begin{equation}
(-1)^{n-1}
\sum_{k=1}^{n-1} 
\sum_{\substack{E \subseteq [k]}} 
\prod_{ j \in E } (x_j-1)
=
(-1)^{n-1}\sum_{k=1}^{n-1} \xvec_{[k]},
\end{equation}
which is exactly the right-hand side in \eqref{eq:derangementSum}.
\end{proof}


\begin{corollary}
For $n,k \geq 1$, we have that
\[
|\{ \pi \in \dArr^e_n : \exc(\pi) = k \}|-|\{ \pi \in \dArr^o_n : \exc(\pi) = k \}| = (-1)^{n-1}.
\]
\end{corollary}
\begin{proof}
 This follows directly by comparing
 coefficients of degree $k$ in \eqref{eq:derangementSum}.
\end{proof}




\subsection{A right-to-left minima analog}

Our goal with this subsection is to prove 
a right-to-left minima analog of \cref{lem:excEvenOdd}.
We shall use the same type of proof, i.e., by exhibiting an involution on $\symS_n$,
such that all fixed-points with the same set of right-to-left minima,
also have the same sign.

\begin{definition}
Let $\kappa : \symS_n \to \symS_n$ be defined as follows.
Given $\pi \in \symS_n$, let $i \in [n]$ be the smallest \emph{odd} integer 
such that $\pi (i \; i+1)$ and $\pi$ have the same sets of right-to-left minima,
if such an $i$ exists. That is, we swap the entries at positions $i$ and $i+1$ in $\pi$.
%
We then let $\defin{\kappa(\pi)} \coloneqq \pi (i \; i+1)$,
and $\kappa(\pi) \coloneqq \pi$ otherwise.
We say that $\pi$ is \defin{decisive}\footnote{As a nod to the word \emph{critical}.}
if it is a fixed-point of $\kappa$.
\end{definition}

\begin{example}
In $\symS_7$, there are $8$ decisive permutations:
\[
1234567,\;
1234657,\;
1243567,\;
1243657,\;
2134567,\;
2134657,\;
2143567,\;
2143657.
\]
Note that $\{1,3,5,7\}$ are always right-to-left minima (but there might be more).
\end{example}

\begin{lemma}\label{lem:ltrMinEvenOdd}
The map $\kappa : \symS_n \to \symS_n$ has the following properties:
\begin{enumerate}[label=({\roman*})]
 \item \label{it:kappaprop1} $\kappa$ is an involution,
 \item \label{it:kappaprop2} $\kappa$ preserves the number of right-to-left minima,
 \item \label{it:kappaprop3} $\kappa$ changes sign of non-fixed elements,
\item \label{it:kappaprop4} 
 For each subset $T \in [n] \cap \{2,4,6,\dotsc\}$,
 there is a unique decisive permutation with 
 $\{1,3,5,\dotsc\} \cup T$ as right-to-left minima set.
  
 \item \label{it:kappaprop5} 
 there are $\binom{ \lfloor n/2 \rfloor }{k - \lceil n/2 \rceil}$
 decisive permutations with exactly $k$ right-to-left minima,
 and they all have sign $(-1)^{n-k}$.
 \end{enumerate}
\end{lemma}
\begin{proof}
Items \ref{it:kappaprop1}--\ref{it:kappaprop3} are clear from the definition of $\kappa$.
%
It remains to prove \ref{it:kappaprop4} and \ref{it:kappaprop5}.
Let us use $O$ to denote the odd integers in $[n]$, and let $E$
be the even integers in $[n]$.
%
In order to prove \ref{it:kappaprop4}, we must construct a decisive permutation $\pi$,
such that $\RLMv(\pi) = O \cup T$. We construct $\pi$ from $T$
subject to the following constraints:
\begin{itemize}
\item if $n$ is odd, then $\pi(n)=n$.
\item if $j \in O$, $j<n$ we have that
\[
\begin{cases}
\pi(j) =j \text{ and } \pi(j+1) = j+1 &\text{ if } j+1 \in T \\
\pi(j) =j+1 \text{ and } \pi(j+1) = j &\text{ if } j+1 \notin T.
\end{cases}
\]
\end{itemize}
In short, $\pi$ is constructed by first placing $1$ and $2$,
in the order determined by $T$, then $3$ and $4$, etc.
By construction, $\RLMv(\pi) = O \cup T$.
Now we must show that a permutation is decisive if and only if they
are in this form. From the construction, it is clear that
$\pi (i\; i+1)$ and $\pi$ will not have the same set of right-to-left minima,
for any choice of $i \in O$. Hence, all permutations with this structure are decisive.
\medskip

\textbf{Claim:} \emph{Every decisive permutation has the structure described above}.
First note that the claim is true for $n=1$ and $n=2$.
Now, if $n$ does not appear among the last two entries of $\pi$,
then $n$ is not a right-to-left minima. Moreover, we can swap this entry 
with the entry either to its right or to its left, and preserve the 
set of right-to-left minima. 
In particular, if $n$ does not appear among the last two positions,
then $\pi$ is not decisive. Now, if $\pi(n)=n$, we can remove 
the last entry and use induction.
Otherwise, suppose $\pi(n-1)=n$ and $\pi(n)<n-1$. In particular,
$n-1 \notin \RLMv(\pi)$.
We can then swap $n-1$ with one of its neighbors (depending on if $n-1$
appears at an odd or even position) and be sure to preserve $\RLMv(\pi)$.
Hence, $\pi$ is not decisive, so we conclude that if 
$\pi(n) \neq n$, then $\pi(n-1)=n$ and $\pi(n)=n-1$ in order for $\pi$
to be decisive. But in this case, we may remove the last two entries,
and proceed by induction. This ends the proof of the claim.
\smallskip 


From \ref{it:kappaprop4}, we know that in order to 
construct a decisive permutations in $\symS_n$ with $k$ right-to-left minima,
we must choose all $\lceil n/2 \rceil$ odd integers in $[n]$,
and pick a subset of size $k  - \lceil n/2 \rceil$,
from the set of even integers in $[n]$.
This subset of even integers has cardinality $\lfloor n/2 \rfloor$.
Hence, we get the advertised formula in 
\ref{it:kappaprop5}, so it suffices to show that all such
decisive permutations have the same sign.
By the claim above, it is evident that 
the sign only depends on the number of right-to-left minima in $\pi$,
and from here, it is straightforward for show that it is indeed
$(-1)^{n-k}$.
\end{proof}


\begin{corollary}\label{cor:rlminsum}
We have that for any $n\geq 1$
\begin{equation}\label{eq:rlminSymSum}
 \sum_{\substack{\pi \in \symS_n   }} (-1)^{\inv(\pi)} \xvec_{\RLMv(\pi)}
 = 
	\big(\prod_{\substack{i \in [n] \\ i \text{ odd}}} x_i\big)
	\big(\prod_{\substack{j \in [n] \\ j \text{ even}}} (x_j-1)\big).
\end{equation}
In particular, for any $k = 1,\dotsc,n$ we have that
\[
 |\{ \pi \in \symS^e_n : \rlm(\pi)=k \}| - 
 |\{ \pi \in \symS^o_n : \rlm(\pi)=k \}| = (-1)^{n-k}
 \binom{ \lfloor n/2 \rfloor }{k - \lceil n/2 \rceil}.
\]
\end{corollary}
\begin{proof}
This follows directly from \cref{lem:ltrMinEvenOdd},
where the first product in \eqref{eq:rlminSymSum}
corresponds to the fact that all odd integers in $[n]$ must be 
right-to-left minima for decisive permutations,
and the second product corresponds to choosing a
subset $T$ among the even numbers in $[n]$.
It remains to check that the signs are chosen correctly,
which is straightforward as well.

The second statement also follows from \cref{lem:ltrMinEvenOdd},
or by simply comparing coefficients of $t^k$ in \eqref{eq:rlminSymSum},
after letting $x_j \to t$.
\end{proof}


We conclude this section with the following problem.
\begin{problem}
Is it possible to formulate an analog of \cref{prop:main}?
In particular, for $T \subseteq [n]$, is there a nice expression for the sum
\[
 \sum_{\substack{\pi \in \symS_n  \\ T \subseteq \FIX(\pi) }} t^{\rlm(\pi)} (-1)^{\inv(\pi)} ?
\]
Computer experiments suggest that this sum is either $0$ or of the form 
$\pm t^a (t+1)^b (t-1)^c$, where $a$, $b$, and $c$ depend on $T$ in some manner.
\end{problem}



\section{Further ideas and conjectures}


\subsection{Multiderangements}

Let $\defin{B_n} \coloneqq (1,1,2,2,3,3,\dotsc,n,n)$ be fixed.
A \defin{biderangement} of $B_n$, is a permutation, $\wvec$, of the entries in $B_n$, 
such that $\wvec(j) \neq B_n(j)$ for all $j \in [2n]$.
The set of biderangements of $B_n$ is denoted $\defin{\bdArr_n}$.
The cardinality of $\bdArr_n$ is given by \oeis{A000459},
which starts as $0, 1, 10, 297, 13756,\dotsc$.

We compute the number of \defin{inversions} 
of a biderangements as for words in general --- see \cref{eq:invDef}.
Moreover, we say that $w(j)$ is an \defin{excedance value} of $\wvec \in \bdArr_n$
if $w(j) > B_n(j)$. This defines the multi-set valued statistic $\EXCv(\wvec)$.
We also define the \defin{right-to-left minima values}
as the set 
\[
   \RLMv(\wvec) \coloneqq \{ w_i : w_i < w_j \text{ for all $j \in \{i,i+1,\dotsc,2n\}$} \}.
\]

\begin{example}
In the table below, we show the ten elements in $\bdArr_3$,
together with the corresponding inversion and excedance statistics.
%
\begin{center}
\begin{tabular}{ccccccccc}
 \text{Biderangement} & \text{inv}  & \text{EXCv} & \text{RLMv} & $\quad$ & \text{Biderangement}  & \text{inv} & \text{EXCv} & \text{RLMv} \\
   \toprule \\
 223311&8&2233&1   &&  231312&7&233&12  \\
 231321&8&233&1   &&  233112&8&233&12  \\
 233121&9&233&1   &&  321312&8&233&12  \\
 321321&9&233&1   &&  323112&9&233&12  \\
 323121&10&233&1  &&  331122&8&33&12 
\end{tabular}
\end{center}
\end{example}


\begin{proposition}
For $n\geq 1$ , we have that 
\begin{equation}\label{bidera}
 \sum_{\wvec \in \bdArr_n} (-1)^{\inv(\wvec)} 
 \xvec_{\EXCv(\wvec)} \yvec_{\RLMv(\wvec)}
 = 
 \sum_{\pi \in \dArr_n} \xvec_{\EXCv^2(\pi)} \xvec_{\RLMv(\pi)},
\end{equation}
where $\EXCv^2(\pi)$ is the multiset obtained from $\EXCv(\pi)$ by repeating each element twice.
\end{proposition}
\begin{proof}
Define a mapping $\beta:\bdArr_n\to\bdArr_n$ as $\beta(w)=w'$, where $w'$ is obtained
from $w$ by switching $w(j)$ and $w(j+1)$ for the smallest odd $j\in[2n-1]$ such that $w(j)\neq w(j+1)$.
If there is no such $j$, then $\beta(w)=w$.
Since $w'(j)=w(j+1)\neq B_n(j+1)=j=B_n(j)$ and $w'(j+1)=w(j)\neq B_n(j)=j=B_n(j+1)$, 
$\beta$ is a sign reversing involution that preserves
the excedance set values, $\EXCv(w)=\EXCv(w')$.
Since the others $w(j)$ and $w(j+1)$ appears to
the right of $j$ and $j+1$, otherwise $j$ would
not be the smallest, the $\RLMv$ will consider
the others $w(j)$ and $w(j+1)$. So $\RLMv(w)=\RLMv(w')$.  
The fixed points are bijection with the derangements,
sending a derangement 
$\pi$ to $v=\pi(1)\pi(1)\pi(2)\pi(2) \dotsb \pi(n)\pi(n)$.
Thus, $\EXCv(v)$ is the multiset obtained from 
$\EXCv(\pi)$ by repeating each element twice 
and $\RLMv(v)=\RLMv(\pi)$ and then equation 
\ref{bidera} satisfied.
\end{proof}



\subsection{A generalization}

The set of dearrangements are permutations where cycles of length $1$
are disallowed. With this in mind, it is reasonable to explore what 
happens if we add restrictions to the length of the cycles.

\begin{conjecture}
Let $k$ be a fixed positive integer.
Then
\[
(-1)^n
 \sum_{\pi \in \symS_n} (-1)^{\length(\type(\pi))}
 \left(\prod_{j\in \RLMv(\pi)} x_j \right)
 \left(\prod_{j\in \EXCv(\pi)} y_{\sigma(j)} \right)
\]
is an element in $\setN[x_1,\dotsc,x_n,y_1,\dotsc,y_n]$.
\end{conjecture}

Note that $(-1)^{\length(\type(\pi))-1} = (-1)^{\inv(\pi)}$,
so the case $k=1$ does indeed correspond to \cref{eq:mainResult}.
For the case $k=n$, all permutations have the same sign, so the statement is trivial.
Interestingly, summing over all permutations consisting of a single cycle of length $n$,
\begin{equation}\label{eq:EulerianSum}
 \sum_{\substack{\pi \in \symS_n \\ \type(\pi) = (n)}}
 \left(\prod_{j\in \RLMv(\pi)} x_j \right)
 \left(\prod_{j\in \EXCv(\pi)} y_{\sigma(j)} \right)
\end{equation}
is a multivariate polynomial where the number of terms (not counting multiplicity!)
seems to be given by the sequence \oeis{A124302}, which starts as
\[
 1, 1, 2, 5, 14, 41, 122, 365, 1094, \dotsc.
\]
For $n=3$, \eqref{eq:EulerianSum} is equal to
\[
 x_1 x_2 x_3 y_4 + 
 x_1 x_3 y_2 y_4 + 
 x_1 y_3 y_4 + 
 2 x_1 x_2 y_3 y_4 +
 x_1 y_2 y_3 y_4
\]
which has $5$ terms.


\subsection{Right to left minima and derangements}

Let $a_{n,k}$ be defined via
\begin{equation}
 \sum_{\pi \in \dArr_n} t^{\rlm(\pi)} = \sum_{k=1}^n a_{n,k} t^k,
\end{equation}
so that $a_{n,k}$ is the number of derangement with exactly $k$ right-to-left minima.
For example, the data for $a_{n,k}$ is shown in the table below.

\begin{table}[!ht]
\begin{tabular}{rrrrrrrr}
 & \textbf{1} & \textbf{2} & \textbf{3} & \textbf{4}& \textbf{5}& \textbf{6}& \textbf{7} \\
\toprule
\textbf{2}& 1 \\
\textbf{3}&1 & 1\\
\textbf{4}&3 & 5 & 1 \\
\textbf{5}&11 & 21 & 11 & 1 \\
\textbf{6}&53 & 113 & 79 & 19 & 1 \\
\textbf{7}&309 & 715 & 589 & 211 & 29 & 1  \\
\textbf{8}&2119 & 5235 & 4835 & 2141 & 461 & 41 & 1 \\
\end{tabular}
\caption{The number of derangements of $[n]$ with exactly $k$ right-to-left minima.}\label{tab:deArrRLM}
\end{table}

It is straightforward from the definition to see that $a_{n,1}$ is the number of derangement ending with a $1$.
The sequence $a_{n,1}$ shows up in the OEIS as \oeis{A000255}, which hints about the recursion
\[
 a_{n,1} = (n-2) \cdot a_{n-1,1} + (n-3) \cdot a_{n-2,1}, \qquad a_{1,1} =1, \qquad a_{2,1} =1.
\]
Moreover, it seems that $a_{n,n-1} = (n-2) + (n-1)^2$.
The data in \cref{tab:deArrRLM} is not in the OEIS.


\subsection*{Acknowledgements}

The authors have benefited greatly from the discussions
with Jörgen Backelin of Stockholm University.
The Online Encyclopedia of Integer sequences \cite{OEIS} 
has been an invaluable resource in our project.

The second author is grateful for the financial support extended
by the cooperation agreement between International Science Program
at Uppsala University and Addis Ababa University.


\bibliographystyle{alphaurl}
% \bibliography{bibliography}

\begin{thebibliography}{Man93}

\bibitem[BV17]{BarilVajnovszki2017}
Jean-Luc Baril and Vincent Vajnovszki.
\newblock A permutation code preserving a double {E}ulerian bistatistic.
\newblock {\em Discrete Applied Mathematics}, 224:9--15, June 2017.
\newblock \href {https://doi.org/10.1016/j.dam.2017.02.014}
  {\path{doi:10.1016/j.dam.2017.02.014}}.

\bibitem[Cha01]{Chapman2001}
Robin Chapman.
\newblock An involution on derangements.
\newblock {\em Discrete Mathematics}, 231(1-3):121--122, March 2001.
\newblock \href {https://doi.org/10.1016/s0012-365x(00)00310-1}
  {\path{doi:10.1016/s0012-365x(00)00310-1}}.

\bibitem[Man93]{Mantaci1993}
Roberto Mantaci.
\newblock Binomial coefficients and anti-exceedances of even permutations: A
  combinatorial proof.
\newblock {\em Journal of Combinatorial Theory, Series A}, 63(2):330--337, July
  1993.
\newblock \href {https://doi.org/10.1016/0097-3165(93)90064-f}
  {\path{doi:10.1016/0097-3165(93)90064-f}}.

\bibitem[MR01]{MantaciRakotondrajao2001}
Roberto Mantaci and Fanja Rakotondrajao.
\newblock A permutations representation that knows what ``{E}ulerian'' means.
\newblock {\em {Discrete Mathematics \& Theoretical Computer Science}}, 4(2),
  January 2001.
\newblock URL: \url{https://dmtcs.episciences.org/271}.

\bibitem[MR03]{MantaciRakotondrajao2003}
Roberto Mantaci and Fanja Rakotondrajao.
\newblock Exceedingly deranging!
\newblock {\em Advances in Applied Mathematics}, 30(1-2):177--188, February
  2003.
\newblock \href {https://doi.org/10.1016/s0196-8858(02)00531-6}
  {\path{doi:10.1016/s0196-8858(02)00531-6}}.

\bibitem[Slo19]{OEIS}
Neil J.~A. Sloane.
\newblock The {O}n-{L}ine {E}ncyclopedia of {I}nteger {S}equences.
\newblock Online, 2019.
\newblock URL: \url{https://oeis.org}.

\end{thebibliography}


\end{document}

